\documentclass[11pt,a4paper]{article}
\usepackage[utf8]{inputenc}
% lmodern must precede fontenc: without a scalable Type1 font the T1 encoding
% falls back to bitmap fonts, and microtype's expansion then aborts the run
% with "auto expansion is only possible with scalable fonts".
\usepackage{lmodern}
\usepackage[T1]{fontenc}
\usepackage[margin=1in]{geometry}
\usepackage{amsmath,amssymb}
\usepackage{graphicx}
\usepackage{microtype}
\usepackage[hidelinks,breaklinks]{hyperref}
\usepackage{url}
\urlstyle{same}

% Generated by code/tools/md_to_tex.py from the SurveyForge Markdown output.
% Citation numbers match the source .md exactly.

\title{Retrieval-Augmented Generation for Large Language Models: Architectures, Methodologies, and Future Directions}
\author{Generated by SurveyForge}
\date{}

\begin{document}
\maketitle
\tableofcontents
\newpage



\section{Introduction}


The rapid advancement of large language models (LLMs) has fundamentally transformed natural language processing, yet these models exhibit critical limitations that constrain their deployment in knowledge-intensive applications. Central among these is the reliance on static parametric knowledge encoded during pre-training, which becomes increasingly obsolete as world knowledge evolves \cite{arxiv:2312.10997v5,arxiv:2410.12837v1}. Empirical evidence demonstrates that LLMs struggle with less popular factual knowledge, and scaling model parameters alone fails to appreciably improve memorization of long-tail facts \cite{arxiv:2212.10511v4}. This architectural constraint manifests most prominently as hallucination---the generation of plausible yet factually incorrect content---and poses significant obstacles to the deployment of LLMs in scenarios demanding temporal freshness, verifiable reasoning, and domain-specific expertise beyond the pre-training corpora \cite{arxiv:2312.10997v5,arxiv:2506.00054v1}. These foundational limitations establish the central motivations for augmenting LLMs with non-parametric, externally sourced knowledge.

Retrieval-augmented generation (RAG) emerges as a paradigm that directly addresses these limitations by integrating external knowledge sources into the generation pipeline \cite{arxiv:2005.11401v4}. Formally, a RAG system couples a parametric generator---typically a pre-trained seq2seq or decoder-only LLM---with a non-parametric memory (e.g., a dense vector index of a knowledge corpus), accessed via a neural retriever \cite{arxiv:2005.11401v4}. By conditioning generation on retrieved evidence, RAG systems augment the LLM's internal knowledge with dynamic, verifiable external information, enabling more specific, diverse, and factual outputs compared to parametric-only baselines \cite{arxiv:2005.11401v4,arxiv:2305.06983v2}. The paradigm has gained significant traction as a solution for grounding responses in up-to-date, domain-specific knowledge repositories, and for providing provenance for model decisions \cite{arxiv:2403.03187v1,arxiv:2405.07437v2}.

Tracing the evolution of RAG reveals a trajectory from early non-parametric and semi-parametric models to modern, tightly integrated modular systems. Foundational work on retrieval-augmented text generation established the basic retrieve-then-generate paradigm, incorporating retrieved memory as guidance for dialogue and generation tasks \cite{arxiv:1809.05296v5,arxiv:2202.01110v2}. The seminal introduction of RAG models, where the parametric memory is pre-trained seq2seq and the non-parametric memory is a dense vector index, significantly advanced the field by unifying the retriever and generator onto a common sequence-to-sequence framework \cite{arxiv:2005.11401v4}. Subsequent evolution witnessed a progression from Naive RAG to modular approaches involving pre- and post-retrieval optimizations, query transformations, and multi-step reasoning \cite{arxiv:2312.10997v5,arxiv:2305.06983v2}, culminating in iterative and agentic architectures that orchestrate multiple components \cite{arxiv:2501.09136v4,arxiv:2411.19443v1}. This trajectory positions RAG not as a singular monolithic solution but as a rapidly diversifying field of methodological approaches.

A comparative analysis with alternative knowledge injection methods is essential to fully understand its design space. Full-model fine-tuning modifies the LLM's parameters to learn tasks but requires task-specific data and risks catastrophic forgetting \cite{arxiv:2402.01364v2}; parameter-efficient methods such as LoRA can mitigate computational demands but still handle knowledge at the parameter level \cite{arxiv:2603.04647v1}. Conversely, long-context window expansion directly supplies raw context to the model, bypassing retrieval but incurring significant computational cost and can degrade performance with a larger retrieved passage set \cite{arxiv:2410.05983v1,arxiv:2407.02485v1}. Knowledge editing offers targeted modifications but does not address the need for continuous updates and large-scale knowledge integration. Within this broad documented spectrum, RAG has demonstrated robust performance for factual tasks, often performing as well or better than other validation-augmented approaches \cite{arxiv:2212.10511v4}; and it can also be combined with other strategies, producing effective hybrid systems \cite{arxiv:2501.15915v1,arxiv:2410.05983v1}.

Within this broad landscape, our survey aims to cut a precise and rigorous path. Rather than compiling an encyclopedic review of every detail, we focus on the integrated RAG system---its core tripartite architecture---encompassing retrieval, augmentation, and generation. This perspective emphasizes the architectural variety, methodological analyses, design decisions, and evaluation and trustworthiness of the underlying integrated systems, distinguishing our contribution from prior surveys that have covered the wider IR or LLM domains \cite{arxiv:2308.07107v3,arxiv:2506.00054v1}. By framing the entire system as the unit of analysis, we seek to provide a coherent and critical examination of how the retrieval component, generation component, and their complex interaction define the capabilities, limitations, and future potential of current RAG systems.


\section{Foundational Architectures and Design Paradigms}



\subsection{Taxonomy of Integration Paradigms}


The architectural design space of retrieval-augmented generation (RAG) systems is fundamentally defined by the nature of the interaction between the retrieval and generation components. This interaction dictates not only the flow of information but also the mechanisms for gradient propagation, the degree of system-level coherence achievable, and the practical constraints on training and deployment. We can formalize a taxonomy of integration paradigms along a spectrum from fully joint, end-to-end trained systems to modular, loosely coupled pipelines, with recent innovations carving out a middle and increasingly influential path through unified, single-model architectures.

At one extreme lie \textbf{pre-training-based integration paradigms}, where the retriever and generator are co-trained from scratch (or substantially fine-tuned together) to share a unified objective. This paradigm, pioneered by works such as REALM and RAG, treats the entire system as a differentiable whole. The retriever is trained via a maximum inner product search (MIPS) loss that is directly supervised by the end-task generator's likelihood, effectively creating a system where the retrieval process is optimized to provide the most useful latent knowledge for generation. As classical work on RAG has shown, this yields a powerful synergy, achieving state-of-the-art results on knowledge-intensive tasks by seamlessly balancing the models parametric memory with non-parametric external knowledge \cite{arxiv:2005.11401v4,arxiv:2312.10997v5}. The primary advantage is a tightly coupled system that can adapt its retrieval strategy for a specific task without human intervention. However, the computational costs and engineering complexity are substantial, and the stability can be a concern, as the retriever's performance can degrade if the generator is not fully trained \cite{arxiv:2506.00054v1}.

At the opposite end of the spectrum lies the \textbf{pipeline or cascade architecture}, which has become the dominant design for practical deployments due to its modularity and efficiency. In this design, a pre-trained, independent retriever and generator are connected via a non-differentiable interface, typically by simply concatenating the retrieved documents to the input prompt \cite{arxiv:2202.01110v2,arxiv:2506.00054v1}. While this separation allows for the independent use of the strongest available models and facilitates human inspection, it creates a profound semantic gap: the retriever is typically optimized for relevance, while the generator is trained for fluency, creating a fundamental disconnect between the two. This divergence is manifested through examples like the poor handling of the "Lost-in-the-Middle" problem by sequence-to-sequence models or the reliance of retrievers on lexical overlap as opposed to the deeper reasoning required by some generation tasks \cite{arxiv:2411.19463v3,arxiv:2410.05983v1,arxiv:2411.19463v3}. To mitigate this, frameworks like DuetRAG and R2AG have been engineered to bridge the gap, but they often require intricate loss functions or additional intermediary modules \cite{arxiv:2507.18910v1}.

Finally, emerging at the state of the art is a \textbf{unified and fully-integrated architecture} that aims to transcend the discord of the previous two paradigms. These systems take a single pre-trained decoder-only language model and divide its layers into a retriever group and a generator group with shared parameters on a unified forward pass. The retriever becomes an internal latent function, rather than an external module requiring a separate dense index. This is perhaps most advanced the notion of parametric RAG, or "RAG as a latent mechanism", elegantly demonstrated under the labels "retrieval as generation" or "Parametric RAG". Here, a model is explicitly trained to generate relevant evidence or "memories" from its weights before generating the final response, thus minimizing the semantic discrepancy between retrieval and language models and nullifying the need for a human-specified explicit query in the system's forward pass \cite{arxiv:2501.15915v1,arxiv:2503.23895v4}. This consolidation can be considered a core step toward building a broad, general-purpose "retrieval-augmented LM" that can perform tasks beyond open-domain QA, positioning retrieval as a valuable feature that can be accessed internally across different domains \cite{arxiv:2403.03187v1}.


\subsection{Retrieval Augmentation Structures and Knowledge Representations}


The architectural efficacy of a Retrieval-Augmented Generation (RAG) system is fundamentally constrained by the nature and organization of its external knowledge source---the substrate upon which the integration paradigms described earlier operate. This subsection dissects the principal knowledge representation paradigms---unstructured text corpora, structured knowledge graphs, and multi-modal/hybrid sources---examining how indexing schemas, granularity, and dynamic update mechanisms collectively determine the upper bounds of achievable grounding and output quality. The choice of representation is not a mere storage decision; it is an architectural commitment that dictates the very semantics of the retrieval process and the subsequent reasoning capacity of the generative model. Indeed, as the architectural designs discussed previously define \emph{how} retrieval and generation interact, the knowledge representation layer defines \emph{what} the system can access from the world.

The most prevalent and foundational structure is the unstructured text corpus, typically partitioned into chunks and indexed within a vector database for dense retrieval. This paradigm, central to classic RAG systems and directly aligned with the pre-training-based and pipeline architectures described earlier, relies on semantic or lexical indexes over raw, vectorized text \cite{arxiv:2005.11401v4}. The granularity of the chunk is a critical parameter, balancing specificity against contextual coherence; overly fine chunks risk fragmenting crucial evidence, while overly coarse chunks introduce noise and dilute relevance \cite{arxiv:2501.07391v1}. To manage the resulting context bloat and improve latency, variants like TurboRAG have explored pre-computed KV caches for chunked text, while cache-augmented generation (CAG) paradigms challenge the need for real-time retrieval by pre-loading full corpora into long-context windows, directly addressing the latency and retrieval-error trade-offs inherent in flat indexes \cite{arxiv:2410.07590v1,arxiv:2412.15605v2}. These optimizations are particularly relevant to the pipeline architecture's emphasis on modular efficiency, as they enable faster integration of external knowledge without altering the core retrieval-generation interface. However, the flat nature of these indexes limits the system's ability to capture inter-document dependencies, leading to fragmented answer generation that fails to address complex, multi-sourced queries \cite{arxiv:2501.13958v3}---a deficit that the following section will address through workflow-level adaptive traversal and logic-aware retrieval.

To overcome the relational deficits of flat representations, structured knowledge graphs have emerged as a powerful alternative, enabling GraphRAG systems to explicitly model relationships and support multi-hop reasoning. These systems construct the knowledge graph with entity relationships, dictating the query-to-graph traversal and pruning of multi-hop paths to reduce search space and redundancy \cite{arxiv:2502.14902v2}. CausalRAG further refines this, extracting causal relationships to preserve contextual reasoning and improve retrieval precision beyond pure semantic similarity \cite{arxiv:2503.19878v3}, fostering the kind of deeper reasoning that the co-trained pipelines struggled to achieve. Furthermore, incremental Graph-RAG frameworks like EraRAG achieve dual scalability---both time and memory---through dynamic index addition without full graph reconstruction, contrasting sharply with the computational overhead of pre-training-based joint systems \cite{arxiv:2506.20963v2} and setting the stage for the adaptive, evolving workflows discussed in the following sections. Additionally, multi-granularity retrieval paradigms such as FunnelRAG introduce coarse-to-fine pipelines that mitigate erroneous negative indexing quality trade-offs, providing a practical balance between computational budget and evidence-gun specificity for task-aware retrieval \cite{arxiv:2410.10293v3}.

As LLM applications broaden, RAG is increasingly required to integrate heterogeneous data modalities from diverse sources, including images, audio, tables, and structured databases. VisRAG demonstrates that processing entire PDF pages as vision-language models proves particularly effective and multi-source to centralized text-extraction errors leads to irreparable loss of positional and overlapped correctness in textual representation \cite{arxiv:2410.10594v2}. To manage heterogeneous sources, research has moved beyond single-agent systems to multi-agent architectures that delegate to specialized agents for relational data, NoSQL stores, or document stores, enhancing concrete query efficiency and modular flexibility \cite{arxiv:2412.05838v1}. This multi-agent coordination naturally connects to the agentic workflow designs highlighted as the current frontier in the following subsection. Similarly, and on a competing axis, parametric-knowledge strategies such as Parametric RAG bypass external storage entirely and inject document knowledge directly into the LLM parameters---enabling deep integration, reduction of online computation, and an alternative to in-context memory access that is only feasible to a limited degree under the static-cache approach of CAG \cite{arxiv:2501.15915v1,arxiv:2510.12668v2}.

In synthesis, the field exhibits a clear evolution from simple static, chunked corpora to dynamic, hierarchical structures that mirror the relational complexity of human knowledge. In tandem with the architectural evolution from externally routed pipelines to latent, single-model designs, the key design decision is no longer just how to retrieve, but how knowledge is stored and accessed---whether in an external store, a continuous cache, or the smoothing and provisioned into the model's own parameter space. The interaction between these knowledge choices and the workflow strategies that follow will ultimately determine the system's ability to scale with growing corpora, maintain synthetic coherence, and deliver robust, grounded outputs in real-world, dynamic environments.


\subsection{Evolutionary System Workflows, and Adaptive Methodologies}


The evolution of RAG system workflows represents a fundamental shift from static, single-pass pipelines to dynamic, feedback-driven architectures capable of complex reasoning and multi-source integration. This progression can be understood along several key dimensions of design freedom: the temporal scheduling of retrieval relative to generation, the presence and nature of feedback loops, and the degree of autonomy granted to the system in orchestrating retrieval actions. Understanding this design space is essential, as the choice of workflow paradigm directly determines a system's capability to handle queries of varying complexity, ranging from simple factoid lookup to multi-hop reasoning tasks \cite{arxiv:2410.12837v1}.

The foundational paradigm remains the single-step retrieve-then-generate architecture, which performs all context retrieval in a single pass before the generation stage. The efficiency and simplicity of this approach, however, is bounded by the model's context window and its susceptibility to distractions from irrelevant retrieved information. This has motivated design variations aimed at streamlining this stage. For example, the "cache-augmented generation" (CAG) paradigm proposes a radical simplification where, for limited and manageable knowledge bases, all relevant resources are preloaded into the model's extended context and their runtime parameters cached, thereby bypassing real-time retrieval to eliminate retrieval latency and potential errors, a particularly valuable proposition in static knowledge-management scenarios \cite{arxiv:2412.15605v2}. Complementing these efforts, hardware and system-level optimizations such as precomputing and storing key-value (KV) caches for document chunks, as exemplified by TurboRAG, have been shown to drastically reduce time-to-first-token (TTFT) latency by eliminating online computation of KV caches during the prefill phase \cite{arxiv:2410.07590v1}.

The next evolutionary stage introduces explicit pre- and post-retrieval optimization pipelines. This stepwise architecture treats data pre-processing and post-retrieval refinement as distinct modules for improving quality. Pre-retrieval work aims to bridge the gap between a user's raw query and the knowledge corpus, while post-retrieval stages refine the retrieved set. However, such pipelines often struggle with complex queries where information is fragmented across documents. To address this, query decomposition and planning have been applied. For instance, the LogicRAG framework dynamically extracts reasoning structures at inference time, decomposing a query into subproblems within a directed acyclic graph (DAG), guiding adaptive retrieval that aligns with query logic without the cost of a pre-built graph \cite{arxiv:2508.06105v2}. Similarly, Graph-based RAG methods, such as PathRAG, attempt to better capture inherent dependencies and structured relationships in text by organizing textual information into an indexing graph and then retrieving key relational paths, using flow-based pruning to reduce redundant information \cite{arxiv:2502.14902v2}.

The most advanced workflows are agentic designs that represent the current frontier. These systems orchestrate multiple tools---search tools, memory storage, APIs---within an iterative loop of actions and feedback, enabling complex problem-solving that explores different paths to a solution. The efficacy of these intelligent workflows has been evaluated in multi-agent environments, and their design has been aligned with theoretical models like Markov Decision Processes to enable sophisticated planning and control \cite{arxiv:2412.05838v1}. The workflow of these systems is often moderated by an evaluator agent to assess the sufficiency of retrieved evidence or final answers, deciding when to stop retrieval or trigger new queries \cite{arxiv:2412.05838v1}. This multi-agent approach is particularly valuable in complex multi-source environments, where specialized agents for distinct data sources (e.g., relational, NoSQL, graph) collaborate to achieve more robust results \cite{arxiv:2412.05838v1}.

Looking forward, the field is moving toward workflows that will not only execute multiple steps but also learn and adapt over time. A central challenge for agentic systems is the significant token consumption innate to multi-turn retrieval. Emerging methods like TeaRAG directly address this trade-off between efficiency and accuracy, compressing retrieved content and reducing reasoning steps via iterative process-aware DPO, achieving significant token reductions while maintaining accuracy \cite{arxiv:2511.05385v2}. Future work will continue to design increasingly dynamic systems that balance power, reasoning effort, and operational cost to meet the varying requirements of complex, real-world tasks.


\subsection{Advanced Interaction: Iterative, Multi-Hop, and Self-Refining Systems}


The evolution of RAG systems from single-shot pipelines toward advanced interaction paradigms addresses a fundamental limitation of static retrieve-then-generate architectures: their inability to decompose complex queries, refine retrieval targets based on intermediate signals, and self-correct in response to knowledge conflicts. Where the previous subsection established how query planning---whether through explicit routing logic in multi-agent systems or compact reasoning structures like LogicRAG's inference-time graph---provides the strategic blueprint for orchestrating retrieval actions, this subsection examines the tactical mechanisms that execute and refine those plans. It focuses on three interlocking families of techniques---multi-step retrieval planning, iterative retrieval--generation feedback loops, and self-reflective knowledge validation---that collectively enable systems to handle tasks demanding the assembly of scattered evidence across multiple documents and reasoning steps.

The progression from predefined decomposition strategies to dynamically self-guiding retrieval policies represents a continuum of increasing adaptability. CoRAG initiates this evolution by formalizing query reformulation as a sequential decision process, training models via rejection sampling to generate intermediate retrieval chains rather than relying on a single retrieval action---a transition that explicitly shifts retrieval from one-time relevance maximization to a trajectory-dependent process conditioned on intermediate generation states. The reliance on compression-based trajectory search, however, limits the scalability and sample diversity of the generated chains, motivating LevelRAG to decouple the planning from the dense retriever-specific optimization by routing each query to the most suitable searcher among sparse, dense, and web-based search modules. Through joint learning of query decomposition and sub-query amplification, LevelRAG plans a multi-step route targeting the most relevant evidence while the retriever itself remains frozen, effectively circumventing the training-propagation bottleneck of CoRAG. The crippling dependence on the iterative retrieval backbone stems directly from the sequential construction of the route, a limitation that GeAR surpasses by augmenting arbitrary base retrievers with graph expansion. By leveraging the graph-specialist model to generate informative multi-hop edges offline, GeAR increases the recall of document connections during the initial retrieval step without curtailing its budget for subsequent iterations. This separation of the structural planning process from the retrieval runtime itself eliminates the need for path-based prediction, consequently adapting the initial retrieval to the multi-hop reasoning requirement.

The adaptive topologies of multi-route planning thus raise a crucial question: how does a system execute sub-problems whose underlying structures are unknown a priori? The tracing of contiguous dependencies enables the iterative retrieval of the most relevant nodes, whereas over-generation risks truncating the search space under limited steps when the target hidden within the interplay of multiple pointers. These very uncertainties invite the complementary mechanisms of adaptive sophisticated conversation, fuzzy search, and knowledge control to resolve the conflict between retrieval sufficiency and distraction avoidance.

\textbf{Iterative retrieval--generation feedback mechanisms.} Beyond pre-planned decomposition, iterative retrieval uses generation outputs to inform subsequent retrieval calls, directly addressing the uniqueness of the queries. Auto-RAG extends this paradigm by logging all retrieved knowledge into an autonomous natural-language memory---the LLM's thoughts are fully expressed in natural language---and consuming fewer tokens per retrieval call. The system then uses a critic to decide whether further retrieval is needed, reducing redundancy in the multi-turn drafting of information-seeking dialogues while the process of moving from LLM-context to context-driven command continues. A dedicated NL-based reasoner with a ProBer that leverages the LLM's hidden states to probe the tokened evidence---unlike "RAG route" requires the construction of knowledge complexity---dynamically guides the retrieval control. The discovery of Ret-Robot, RRAG, and IC-RAG models can be taken as follow-up consolidations, refining latent knowledge plausibility-based pathways by applying early knowledge alignment, positioning a corpus-aware feedback mechanism that avoids unnecessary state-space expansion and improves retrieval precision across diverse target documents. Importantly, the whole framework is optimized end-to-end via a trajectory-value-aware synthesis---generating diverse successful conversation paths and failures under strategy-mixed self-practice to ensure high data efficiency, and secondly, a critic trained on those paths provides a local-scope dynamic checkpoint as the only stop signal, avoiding the user's attention to the detected level of unrelated retrieved distractors in the inference loop.

\textbf{Self-reflection, knowledge validation, and model-state self-correction.} A critical challenge embedded in interacting systems is resolving conflicts between retrieved contexts and the model's own parametric knowledge. These conflicts, especially with knowledge-inconsistent data, distort retrieved document ordering and lead to conceptual misalignment. First-layer targeted ideas act at the source level, inflating the propagation-worthy, logically necessary information of each triple unit. The second-level uses human-like reasoning to question " if the chain breaks when trying to disconnect any two criteria " from a proof-verification standpoint, yielding the justification targets that ultimately make the answer more accurate \cite{arxiv:2604.12610v1}, thereby separating proofs from dependencies and detecting whether a given retrieved document is logically sufficient for the current query. Once generated, iteratively verifying one's own answer in multi-round settings also avoids amplifying retrieved noise in follow-ups when the state space is closed. For practical efficiency, RbFT \cite{arxiv:2501.18365v1} instead injects controlled retrieval defects during generation fine-tuning, teaching the model to be resilient to differences in the retriever's behavior and maintain performance when retrieval quality degrades. Minimal-cost filters, meanwhile, propose gating retrieval on uncertainty: when inherent knowledge or the retrieved corpus is assessed as sufficient/low-confidence, the system skips unnecessary context injection \cite{arxiv:2410.13339v2}. Several approaches extend the conflict-checking recursive perspective, such as the mechanism of translating downstream consumer logic into dependence-aware contrast distillation to adaptively switch off the retrieval according to the competition of its own knowledge transmission.

Taken together, the techniques discussed in this subsection show that the movement from single-step retrieval to multi-hop, self-refining systems can be understood as increasing the number of times the information loop process intermediate representations---on the one hand through query planning, and on the other through generative retrieval refinement at various levels. A central observation is that explicit feedback or explicit state-level connection to LLM hidden states must be managed so that no unconstrained efficiency loss occurs \cite{arxiv:2511.05385v2}, as many multi-step methods today can become bottlenecked by the additional token generation or inference steps. The external knowledge layer detailed in the following subsection will be intimately tied to these interaction patterns, since both multi-hop path traversal and dynamic verification loops fundamentally hinge on the memory layout and indexing granularity required to sustain them---a flat vector store suffers from rapid degradation in noisy contexts, while a hierarchical index naturally supports the state transitions inherent in an iterative cycle. Future work linking the two thrusts should jointly optimize the repository scheme and retrieval state so that the dynamic operations from this subsection can be satisfied, the retrieval graph can be updated continuously with minimal forgetting \cite{arxiv:2512.20144v2}. An opportunity equally critical is the explicit feedback links between self-reflective knowledge validation and the federated routing of the corpora, as the prediction of an agent must be validated from source diversity but also grounded in inference, without cascading error, as evidenced by the limits verified in planning performance.


\subsection{Knowledge Base Representation, Indexing, and Dynamic Updating Systems}


The representation, indexing, and lifecycle management of external knowledge constitute the foundational substrate upon which the entire retrieval-augmented generation (RAG) pipeline operates. These design choices not only determine the theoretical ceiling of retrieval quality but also define the system's operational boundaries in terms of scalability, latency, and adaptability. The evolution from static, analysis-oriented corpora to dynamic, heterogeneous knowledge infrastructures represents an imperative architectural shift, as the assumption of a constant knowledge base no longer holds in real-world deployment scenarios where information evolves continuously. A critical first consideration is the preprocessing pipeline that transforms raw heterogeneous documents into queryable representations. For unstructured text, chunking strategies and embedding schemas must be designed with an acute awareness of the information granularity required by downstream tasks. The hierarchical nature of many real-world documents---such as books or technical manuals---affords a natural logical structure that can be exploited via specialized index structures like BookIndex, which extracts a hierarchical table-of-contents representation and maps entity relationships to tree nodes for more precise query routing \cite{arxiv:2512.03413v1}. Beyond raw chunking, the transformation of unstructured text into structured knowledge graphs (KGs) has emerged as a powerful representation strategy. Systems such as KG-Retriever leverage a hierarchical index composed of a knowledge graph layer and a collaborative document layer, capitalizing on the associative nature of graph structures to strengthen intra-document and inter-document connectivity and mitigate the information fragmentation problem common to multi-hop reasoning \cite{arxiv:2412.05547v2}. This shift toward graph-based indexing is not merely a representational change but an architectural decision that enables path-based reasoning and navigation, offering a more epistemically sound substrate for complex queries compared to flat vector spaces.

The choice between dense vector indexes, sparse inverted indexes, and graph structures is rarely binary; practical deployments are increasingly moving toward hybrid representations. The FRAG framework exemplifies this flexibility by estimating a query's hop range and classifying it as simple or complex, thereby routing different queries to tailored pipelines---vector-based retrieval for straightforward questions and iterative multi-hop KG traversal for complex reasoning \cite{arxiv:2501.09957v2}. Similarly, LogicRAG abstracts away from pre-built static graphs altogether, using the language model to dynamically construct a logic-aware directed acyclic graph of subproblems at inference time, effectively rendering the "index" a byproduct of query understanding rather than an offline corpus artifact \cite{arxiv:2508.06105v2}. This paradigm shift---from static storage to dynamic structure---highlights a key tension in knowledge base design: the utility of an index must be balanced against the cost of its construction, maintenance, and update latency. In this context, the dynamic addressing of RAG as a paradigm involves not just \emph{how} a query is answered, but how the corpus itself changes over time.

The management of dynamic knowledge---where data evolves as a continuous stream---presents a distinct set of computational and architectural challenges. Rather than invoking costly, global re-indexing to capture updates, contemporary approaches favor, and discuss, albeit cautiously, incremental update mechanisms that operate on a unit of new information to insert or refresh local index segments without recomputing the entire graph structure. A critical avenue in dynamic environments is the establishment of adaptive addressing schemes that "commit" to relevance over time, placing the computational burden of retrieval on the query side as new data is observed. Furthermore, the corpora are rarely internally consistent: they are drawn from broad, disparate, and web-accessible sources that are often perpetually evolving \cite{arxiv:2411.11895v1}. In these cases, the RAG system must have mechanisms to gracefully handle knowledge that has become obsolete or has been superseded---the so-called anti-knowledge problem---where the retrieved context directly conflicts with updated ground truth. This makes the static models brittle and necessitates the use of dedicated routing and validation layers to adjudicate between the retrieved content and the parametric knowledge of the LLM \cite{arxiv:2505.23794v2}. The necessary discovery is that in the emergent era of massive web-scale data, retrieval systems must not be conceived of as "pure" search engines, but as complex, distributed data infrastructure with intelligent integration layers, akin to federated search systems orchestrating efficient routing to multiple heterogeneous sources \cite{arxiv:2502.19280v2}. Effective RAG in the future will require its knowledge layer to be an \emph{active} management system, capable of orchestrating the most optimal source---be it a local vector store, a structured KG, or a live Web API---and understanding query-dependent data taxonomies, with data freshness and provenance becoming primary index constraints for diverse, real world application domains.


\section{The Retrieval Component: Methods and Advanced Strategies}



\subsection{Query Transformation, Augmentation, and Reformulation}


The efficacy of a retrieval-augmented generation (RAG) system is fundamentally bounded by the quality of the initial retrieval step, which in turn hinges on the alignment between the user's query and the indexed representation of the knowledge corpus. Raw user queries, often characterized by ambiguity, underspecification, and lexical mismatch with the target documents, constitute a significant bottleneck. Query transformation, augmentation, and reformulation directly address this bottleneck by adapting the input query to better navigate the semantic and lexical space of the knowledge base, thereby expanding initial recall and precision \cite{arxiv:2202.01110v2}. The design of these pre-retrieval techniques is a deliberate strategy to optimize the information-theoretic interface between a typically short, information-sparse user utterance and a vast, heterogeneous document collection. This stage is often the most cost-efficient point of intervention in the entire RAG pipeline, as iterative refinement here precludes downstream generation failures.

A primary family of techniques centers on query rewriting and expansion, which seek to reconcile the query's surface form with the rich, often implicit, vocabulary of the target corpus. This is particularly crucial for lexical retrieval models like BM25, which operate under a strict term-matching paradigm and are heavily penalized by vocabulary mismatch. In this context, generating multiple paraphrased variants of the original query \cite{arxiv:2501.07391v1} provides a simple yet effective ensemble approach, maximizing recall by covering a broader lexical surface area. More sophisticated methods, such as query-aware rewriting, leverage in-context learning to elicit a more precise, likely answerable response, with the improved query often incorporating domain-specific terminology or explicit fact-relations absent from the original phrasing \cite{arxiv:2501.07391v1}. Beyond single-turn reformulations, models are also being used to decompose a target task into sequentially searchable sub-questions, each of which is searched independently against the source of evidence \cite{arxiv:2502.11371v3}. The strategic choice between producing a single, analogously specific expansion versus multiple parallel fin-grained sub-queries is a key design decision, depending on the complexity and multi-faceted nature of the query and the downstream task.

Another frontier in this space is breaking down the temporal and semantic distance between the query and the corpus. Pseudo-document generation moves beyond surface-level rewrites by leveraging the LLM's generative capabilities to create a hypothetical document or answer that explicitly verbalizes the query's underlying intent. This generated answer serves as the final retrieval query, operating under the hypothesis that it lexically resembles the gold documents to be retrieved more closely than the original query. There are two primary paradigms for this: the first relies on zero-shot prompting, either through a single forward pass or in a Chain-of-Thought (CoT) manner where model partially reasons before outputting the pseudo-answer; the second, more advanced paradigm, treats pseudo-answer generation as a reinforcement learning problem \cite{arxiv:2501.18365v1}, optimizing the LLM's decoding strategy to produce a pseudo-document specifically designed to maximize the downstream retrieval recall. While the zero-shot approach is extremely flexible and cheap, the RL-optimized generation shows notable performance gains on fact-heavy tasks given its ability to escape myopic behavior, albeit at a computationally cost during the generation phase.

Beyond the explicit generation of text, an emerging category---implicit and query-free retrieval---completely abstracts away from learning the surface form. Instead of generating an explicit query as an intermediate artifact, these systems directly predict a retrieval vector from the LLM's internal hidden states, using a separate attention mechanism to decode it from the context. This setup significantly reduces the complexity of decoupling retrieval from generation, yet it may also limit the expressive power of the query if not carefully trained. \cite{arxiv:2005.11401v4,arxiv:2405.06211v3}. The evolution and the design space of the pre-retrieval stage highlight the shift from improving the quality of the query for a static retriever towards jointly addressing the query and the downstream generation objective, holding the balance between per-query accuracy and whole-pipeline efficiency.

The evolution of query processing in RAG systems, particularly with agentic and iterative retrieval paradigms, has transformed the query from a static input into a dynamic process. Iterative-querying systems, such as FLARE \cite{arxiv:2305.06983v2} and Auto-RAG \cite{arxiv:2411.19443v1}, now treat query transformation as an adaptive process conditioned on the results of the top-k retrieval, allowing for iterative cycles of refining both the query's form and the retrieval space. A comprehensive accounting of future demands reveals that the query process must be unified with the \textbf{generation} process to handle multi-step reasoning tasks, highlighting the need for tighter coupling between the query form and the final generation objective. The cost of excessive reliance on query-based, step-by-step refinement is an increase in the fact that the computational overhead and internal knowledge, such that a "perfectly formed" query is in direct competition with an under-specified, open-ended query guide to handle incomplete retrieval paths. The new frontier lies in building truly modular, hybrid frameworks---the so-called "infrastructure for scaling"---that seamlessly combine query transformation with the dynamic retrieval and generation steps \cite{arxiv:2310.11511v1,arxiv:2410.12837v1}.


\subsection{Taxonomy of Retrieval Models and Algorithms}


The retrieval algorithms that index and search external knowledge form the computational backbone of any RAG system, and the choice of retrieval model establishes the upper bound on the quality of the entire pipeline. The taxonomy of these methods spans a spectrum from classical sparse lexical systems to modern dense, hybrid, and graph-based approaches, with each paradigm offering a distinct trade-off between efficiency, effectiveness, and scalability that practitioners must carefully navigate. Understanding this landscape is essential, as it directly sets the conditions for the query transformation techniques discussed in the previous subsection---the sophistication of query reformulation is inextricably linked to the retrieval model's capacity to interpret the transformed query---and simultaneously defines the quality of the raw candidate list that the post-retrieval refinement stage, described in the following subsection, must contend with.

Sparse retrieval models, exemplified by TF-IDF and BM25, represent the foundational class of lexical matching algorithms. These systems index documents by term frequency-inverse document frequency (TF-IDF) statistics, which can be formalized as follows: given a query \(q\) and document \(d\), BM25 scores are computed as a linear sum of term-weighting functions based on term frequency saturation and document length normalization. These models exhibit robustness and efficiency due to their use of sparse, inverted indexing structures, which allow for sub-linear query time with respect to corpus size. Recent large-scale analyses position lexical retrieval as a strong baseline that often translates well to modern RAG systems without fine-tuning \cite{arxiv:2508.06401v3}. Furthermore, scaling studies investigating the balance between tokens used for pretraining versus retrieval demonstrates a superlinear relationship with document count, highlighting that BM25 and its extensions can scale to billion-scale corpora \cite{arxiv:2604.00715v1}. However, despite these strengths, lexical methods inherently suffer from the lexical gap problem, failing to match queries and documents that are semantically related but share no token overlap. Given that the previous subsection shows how query rewriting and expansion can produce lexically diverse variants, the consistent absence of semantic sensitivity remains the primary limiting factor that dense retrieval models arose to conquer.

Dense retrieval models using bi-encoder architectures---such as DPR---use two distinct neural encoders to project the query and document into a shared low-dimensional embedding space and compute similarity via inner product of the representing vectors. These models not only bridge the semantic gap left to be exploited by sparse methods, but also align naturally with the pseudo-document generators in the pre-retrieval stage, since both strategies operate from a representation determined by learned semantic density. These models leverage Transformer representations and can be fine-tuned with supervision, bridging the performance gap from both lexical and even in few-shot settings \cite{arxiv:2401.14887v3}. While these dual-encoder architectures excel at capturing the conceptual content of text, they generally require dedicated vector bottlenecks. They create a dependency on approximate nearest neighbor search in large-scale vector indices with memory and software overhead, and, unlike lexical methods, retrieval quality may be reduced for domain-specific out-of-distribution data requiring significant fine-tuning \cite{arxiv:2312.05934v3}.

Given the complementary for different viewing angles on retrieval, the necessity for a tight integration becomes clear, and they frequently pair well with the pre-retrieval rephrasing techniques described above. A popular method is hybrid retrieval, which computes a shared score via linear interpolation of lexical score and dense vector component: where the lexical component provides deterministic matches for domain-specific technical terms while the dense one ensures coverage of beyond-synonym diversity. Some modern systems, however, explore the more direct contribution of LLMs to redefine retrieval, either through prompt-guided active traversal or by leveraging relational reasoning over knowledge graphs. For instance, constructing indexes based on graph structures---as in GraphRAG---efficiently represents semantic relationships in an external corpus, allowing retrieval by navigating relational paths via graph traversal \cite{arxiv:2502.14902v2}. Similarly, retailer systems such as LightRAG ensure optimal multimodal retrieval via graph-based learning \cite{arxiv:2410.05779v3}, whereas other graph-centric methods use textual association chains to enable associative discovery and deeper inter-relation highlights of evidence \cite{arxiv:2604.10426v1}. This evolution increasingly blurs the distinction between the index as a static structure and the retrieval step as a reasoning component, instigating direct synergy with the agentic refinement directions discussed later.

In summary, the retrieval model forms both the practical upper bound and the architectural template that queries must satisfy. The optimal choice depends on the balance of multiple trade-offs, but the patterns observed in modern systems move clearly to retriever either ignoring or explicitly taking advantage of the respective lifetimes of lexical determinism: intentionally covering the system with diversity while being supplemented by the efficient aggregation of graph structure. Empirical evidence indicates that retrieval performance rises substantially with data scale, while knowledge allocation between pretraining and retrieval expansion remains an open area of active research \cite{arxiv:2604.00715v1}. ESSENTIAL---while the challenges of carefully balancing guided retrieval strategies are evident \cite{arxiv:2412.15529v4}, the pathway from purely static index to an adaptable, interactive, and even a modelic interface---one where LLMs actively guide the traversal of retrieval---promises to change benchmarks, making the retrieval stage a coequal component in the generation pipeline. The best choice of retriever, therefore, is the one that interfaces seamlessly with the query transformation that precedes it and lays the coherent evidence block that the later refinement will threshold and compress. This choice ultimately dictates the optimism of the entire architecture.


\subsection{Post-Retrieval Refinement and Re-ranking}


The retrieved document list, while often high-recall, is rarely directly consumable by a generative model; it is characterized by redundant passages, noisy or off-topic content, and a ranking order that optimizes retrieval scores rather than downstream generation utility. Post-retrieval refinement and re-ranking constitutes the critical interface where this raw set is transformed into a compact, ordered, and maximally informative context. This stage, which is often glossed over in simpler architectural descriptions, is a primary determinant of final system quality, for it directly governs the quantity and quality of evidence made available to the parametric generator.

Research demonstrates that the retrieved context's composition is often more consequential than the raw retrieval algorithm itself, with poor refinement inherently limiting the potential of even the best retriever \cite{arxiv:2404.10981v1}. Refinement functions can be formally decomposed into two related operations: relevance re-ranking, which re-orders the corpus to attend the most useful evidence, and information compression, which aggregates multiple chunks into a denser, less redundant representation. These operations can be applied independently or jointly but each has distinct technical implications with respect to inference cost, performance, and semantic fidelity.

\textbf{Relevance Re-Ranking and Lightweight Filtering.} The most established class of post-retrieval refinement methods focuses on re-scoring the initially retrieved candidates with a more sophisticated relevance model than the first-stage retriever. The workhorse of this phase is the cross-encoder, a model that jointly encodes the query alongside a candidate document, yielding a more sensitive interaction-level relevance score than the bi-encoder pre-ranking. While notably more expensive, the cross-encoder offers an important trade-off---they can be swapped for lighter-ranking models, such as pointwise or pairwise rankers, to balance computational cost and accuracy. In large-scale deployments, cross-encoders do not scale well due to the cost of encoding numerous passages in a short time, so many systems opt for more lightweight filters. These methods use statistical signatures, like information entropy or token-level redundancy, to dynamically filter out "low-quality" documents without independent deep-language understanding \cite{arxiv:2504.19754v1}. This yields a pronounced practical implication: the use of lightweight filters coupled with re-ranking, i.e., the method can discard noisy context cheaply before a more expensive encoders is applied, is becoming a standard industrial best practice.

\textbf{LLM-Based Listwise and Agentic Re-Ranking.} The boundary between re-ranking and generation has started to blur with the advent of large language models as proxies for re-ranking functions. Unlike cross-encoders, which rely on bi-layer function, LLMs can assess a document from a more holistic, query-relevant semantics, capturing interplay between query and document in a context-aware manner that is often missed by a simple relevance index. Systems increasingly prompt an LLM to function as a list-wise ranker, evaluating and ordering a set of candidate documents in a single forward pass as a unified ranking computation.

This approach allows the ranker to conduct a holistic comparison on the granularity of a set, which can avoid intra-document adhesions that pointwise systems fail to model, and identify more complex information need satisfaction. Other approaches seek to cut computational overhead by preserving the entire list in the prompt both the model to collectively reason over all candidates while minimizing the number of API calls. More recently, this is being factored through the lens of reinforcement learning. For example, systems such as \emph{TeaRAG} incorporate a knowledge graph, which captures semantic similarity and co-occurrence between chunks, and use a joint reward that promotes both answer accuracy and token efficiency, resulting in a re-ranker that learns to prioritize the most \emph{generatively useful} content---thereby directly optimizing the downstream token economy \cite{arxiv:2511.05385v2}. It is a significant stepping stone, as it optimizes ranking for the purpose that is the most rate-determining, i.e., lower tokens---and with no sacrificing output quality.

\textbf{Information Compression and Context Reconstruction.} In the face of long documents or numerous retrieved passages, it is often not only the relevance that matters, but also the sheer length of context that degrades effectiveness. Information compression, or context reconstruction, targets the embedded redundancy and overlapping information. The underlying goal is not to delete relevant documents per se, but to yield a more compact, weighted representation that is focused on the query. Such strategies include query-aware summarization: models can parse the retrieved chunks in their raw form and generate focused summaries before feeding them to the generator, thereby reducing noise, removing unnecessary details that do not relate to the query, and, in some designs, aggregating multiple chunks into a shared synthesized context that incorporates the semantics of all retrieved evidence.

This becomes crucial to address the observed "lost-in-the-middle" problem, which is partially an intrinsic issue with the LLM's attention allocation, but is substantially influenced by the order and information density of the context; correct ordering and sufficient compression largely can alleviate, but not fully solve, the issue. Overall, the field is moving rapidly toward \emph{unified} and \emph{generatively-aware} refinement. While the classical cross-encoder remains a solid baseline, the future of this area lies in a deeper synergy between retriever, reranker, and the generative model itself. The field will see systems moving away from static, one-off ranking, to \emph{multi-functional} and \emph{dynamic} modules that interact with the LLM in a feedback loop, where the LLM can both cause re-ranking iterations and consume the final compact context in a single embedding space.

As RAG frameworks continue to scale their knowledge stores, an increasing differentiating factor is not merely the number of retrieved documents but the ability to filter, optimize, and compress them into a minimal, focused context that supports accurate and robust answer generation. The central remaining challenge is task-aware and cost-aware context compression, which requires fine-grained reasoning over the input, balancing the utility of the accumulated evidence against the finite context budget and the LLM's active capabilities.


\subsection{Adaptive and Dynamic Retrieval Mechanisms}


The necessity of retrieval is not uniform across queries; it depends intrinsically on the boundary between a language model's parametric knowledge and the demands of a specific input. While the previous section established that post-retrieval refinement is critical for shaping raw evidence into an optimal context, an equally pivotal question precedes that stage: \emph{should retrieval even occur, and if so, to what extent?} A growing body of work has therefore moved beyond static retrieve-then-generate pipelines toward adaptive mechanisms that use internal signals or agentic planning to determine if, when, and how to query external sources. This paradigm is driven by two central considerations: the significant computational cost of external memory access and the empirically documented risk that injecting irrelevant or noisy context can degrade generation quality, sometimes even below the performance of a no-retrieval baseline \cite{arxiv:2401.14887v3}. The central challenge is to formally define and measure the model's ``knowledge sufficiency'' at inference time, a problem which has led to a spectrum of strategies ranging from lightweight gating classifiers to sophisticated reinforcement-learning-based policies.

One prominent category relies on confidence signals to gate access, often providing a middle stratum between unconditional retrieval and no-retrieval to optimize multi-turn interaction. Adaptive-RAG introduces a complexity-aware classifier---a smaller language model trained to predict query complexity---that routes between no-retrieval, single-step retrieval, and multi-step iterative approaches. The system balances accuracy against computational overhead by selecting plausible inference pathways \cite{arxiv:2403.14403v2}. Building on this premise, MBA-RAG reframes the routing problem as a multi-armed bandit, where each arm represents a retrieval strategy and an agent learns to dynamically select the optimal action based on query properties. This approach captures the intrinsic exploration-exploitation trade-off in strategy selection and combines multiple single-hop and multi-hop datasets to achieve state-of-the-art accuracy while reducing retrieval cost \cite{arxiv:2412.01572v4}.

Probing internal model states is another effective approach for the ``should I retrieve'' decision, detecting hidden insufficiency indicators rather than relying solely on greedy, top-1 lexical scores. Probing-RAG uses a pre-trained prober that analyzes intermediate hidden state representations to assess the model's internal knowledge state and adaptively determines whether further retrieval is necessary. By summarizing token-level representations of the model's emergent knowledge, it gates access to the retriever only when internal uncertainty is high, reducing redundant steps across open-domain QA benchmarks \cite{arxiv:2410.13339v2}. Similarly, Skill-RAG reveals the typed nature of post-retrieval failures, integrating a lightweight prober into a router to produce failure-state diagnostics. It detects query-evidence misalignment with the prober and then uses a skill router to modify retrieval actions (e.g., defining terms, generating sub-questions), and improves performance on hard cases without altering model weights \cite{arxiv:2604.15771v3}.

In the agentic paradigm, LLMs themselves act as the primary driver of retrieval strategy, and reinforcement learning provides a principled framework for deciding on retrieval actions to maximize long-term value. s3 proposes a searcher-generator protocol that decouples retrieval from generation, training a lightweight search agent using a gain-based reward that encourages retrieval only when it materially improves the probability of a correct response. This search strategy is model-agnostic, and with a small number of curated training outcomes it yields substantial gains on QA benchmarks \cite{arxiv:2505.14146v2}. In another instance, SIM-RAG renders this process as a two-stage framework, training a critic that evaluates whether the current information is sufficient for a correct answer. If insufficiency is detected, the model continues to collect evidence; otherwise, an actor performs the final generation, extending adaptive retrieval beyond single-pass gating to multi-turn iterative loops \cite{arxiv:2505.02811v2}. At a more macro level, agentic retrieval architectures (e.g., hierarchical frameworks) challenge the classic one-shot pipeline and abstract common query-loop patterns into structured components---keyword search, semantic retrieval, iterative question decomposition---that the LLM calls adaptively depending on the information needs of the task \cite{arxiv:2501.09136v4} \cite{arxiv:2602.03442v1}.

Despite the diversity in decision signals, these mechanisms face common trade-offs. The performance of classifier-based gating depends on the estimated query complexity or recall sufficiency, both of which can be fragile in out-of-distribution scenarios where the detector may not generalize. Meanwhile, agentic and reinforcement-learning approaches often incur higher upfront training costs and require reliable reward signals, but they offer more fine-grained control over the cost-accuracy frontier. These methods also raise the question of how to best measure ``dynamism''---for instance, whether to use the number of retrieved rounds, the diversity of retrieved sources, or token-level efficiency as a proxy for retrieval value. Even though techniques like s3 and SIM-RAG demonstrate that a smaller, lighter retrieval policy can outperform static dense pipelines, they also show that no single retrieval strategy is universally optimal.

Looking forward, the field is moving toward more expressive and learnable control policies---via reinforcement learning and RL-distilled heuristics---rather than manual triggers. Integrating these adaptive decisions more closely with generative modeling is a key direction: making the gating decision part of the same architectural flow as the generation model itself, such that retrieval becomes a soft, attuned component of a unified hidden-state process. Another promising avenue is joint optimization of adaptive retrieval and refinement, where the controller not only decides \emph{whether} to retrieve, but also actively influences the construction of the context (e.g., using structured knowledge representations like triplets or graphs) to reduce token overhead and improve downstream grounding \cite{arxiv:2511.05385v2}. Such synergy between adaptive selection and efficient refinement will be essential for scaling RAG to long-horizon, production-ready applications where computational reliability and contextual fidelity must be balanced in real time---paving the way for systems that not only know \emph{what} to retrieve but also \emph{when} to stop searching and focus on generating.


\section{Generation Paradigms and Knowledge Integration}



\subsection{Context Integration and Evidential Prompt Engineering}


The integration of retrieved evidence into the generation prompt is far from a neutral act; the manner in which documents are formatted, ordered, and framed constitutes a critical interface between non-parametric knowledge and the parametric machinery of large language models (LLMs). This stage, often termed evidential prompt engineering, directly modulates the model's attention allocation and its capacity to discriminate between useful and distracting information. The choice of prompt architecture is thus a powerful, low-cost lever for controlling output faithfulness and quality, yet its underlying mechanics are only beginning to be understood.

The most thoroughly documented phenomenon governing context integration is the "lost-in-the-middle" effect, where LLMs exhibit a pronounced performance degradation when salient information is located in the middle of a long input context \cite{arxiv:2410.05983v1}. This U-shaped attention pattern has significant implications for RAG, as retrieved passages are often ranked by relevance and appended to the prompt in descending order of importance. This naive ordering places the most crucial evidence at the periphery---which models process most effectively---and the lower-ranked but potentially necessary information in the attention-poor middle. Prior work demonstrated that re-ordering retrieved passages to present the most relevant ones at the beginning and end of the context is a simple yet powerful training-free optimization, yielding substantial performance gains \cite{arxiv:2410.05983v1}. However, they also find that this simple reordering can be insufficient when many relevant documents are included; a "hard negative" in the middle can still degrade performance. This finding underscores the necessity of pairing reordering with other strategies that enhance the robustness of the generator \cite{arxiv:2410.05983v1}. This interaction between prompt-level engineering and model-level training suggests that the optimal context format is not an intrinsic property of the text but is co-determined by the LLM's own processing dynamics and training objective.

To mitigate positional biases and create a more robust context, various structural justifications have been proposed to highlight evidence. This can take the form of explicit syntax to structure the evidence. Coordinated Semantic Alignment and Evidence Constraints for Retrieval-Augmented Generation with Large Language Models explicitly proposes imposing an evidence constraint \cite{arxiv:2603.04647v1}. This method transforms retrieved context into a core control factor during generation, rather than being interpreted as an implicit context token. This reduction of noise and explicit constraint of content has a clear effect, leading to stable improvements in factuality \cite{arxiv:2603.04647v1}. The framing of evidence as an authoritative anchor also has implications for the handling of knowledge conflicts, i.e., tension between non-parametric retrieval and internal, parametric knowledge. A set of methods analyzes retrieval imperfections and knowledge conflicts and uses an iterative and source-aware prompting strategy that filters and mixes internal and external knowledge with a clear and decomposable formulation \cite{arxiv:2410.07176v2}. Their method adaptively elicits internal knowledge to assess and enhance the final answer, showing the need for the prompt to guide the model in negotiating between sources and potentially overruling the external context \cite{arxiv:2410.07176v2}. This means that instruction-based prompting can act as an extra safety barrier and preference-constructor for the information which the system is about to generate, mitigating hallucinations.

The critiques of single-pass context integration have generated a class of more dynamic, interactive prompting techniques. Rather than processing all retrieved content in a fixed order, these involve the LLM in an iterative process of querying or summarizing evidence, thereby acting as a double-check on the most critical parts of the original context to answer \cite{arxiv:2410.19572v5} or synthesizing new instructions that are a mixture of retrieved and parametric knowledge \cite{arxiv:1809.05296v5}. These methods point towards a paradigm where "context engineering" is not a one-time, deterministic step, but an adaptive agent with the model itself. As future RAG systems become more agentic, we should have a more adaptable, dynamic and evidential prompt control system, activating it on the fly based on the internal state and the confidence of the query itself, making the choice of whether and how to use the retrieved content part of the generative planning process.


\subsection{Parameter-Efficient and Fine-Tuning Strategies for Grounded Generation}


Training-time adaptation of the generator is critical for bridging the gap between generic language modeling objectives and the demands of grounded generation, where outputs must not only be fluent but also faithful to, and appropriately critical of, retrieved evidence. This subsection examines the spectrum of training-time interventions, from full-model fine-tuning to parameter-efficient methods and advanced alignment, with a particular focus on how these strategies shape the model's ability to selectively utilize evidence, reconcile conflicting sources, and balance parametric knowledge with non-parametric context. While the previous subsection highlighted prompt-level strategies to shape context integration at inference, this subsection shifts focus to a more fundamental and durable layer of control: reshaping the model's own parameters to internalize an optimal evidence-utilization policy. These training-time adaptations ultimately determine the upper bound of what prompt engineering can achieve at test time, since the model's learned predispositions---its baseline susceptibility to positional bias, its tendency to over-rely on parametric priors, and its sensitivity to conflicting evidence---directly condition the effectiveness of any evidential prompt construction. Moreover, the robustness instilled through these training-time interventions sets the stage for the lighter-weight, decoding-time corrections covered in the following subsection, which operate on top of the base model to further enforce evidence faithfulness.

During standard pre-training, LLMs internalize a vast but static parametric knowledge base. When augmented with retrieval, the model is confronted with an external, dynamic, and often redundant source of information. The goal of fine-tuning for grounded generation is to teach the model a policy for evidence utilization---one that can distinguish between corroborating and contradictory information. The definitive approach to achieving this is through supervised instruction-tuning on curated query-context-answer triples. This process employs task-specific instruction data to teach the model to leverage the provided context as an authoritative reference. The substantial cost of collecting such annotations has driven the development of methods that leverage synthetic data generation \cite{arxiv:2306.05212v1}. Furthermore, this fine-tuning should also induce robustness to context noise, thereby making the model less brittle to the inevitable imperfections of real-world retrieval output, complementing the prompt-level denoising and reordering strategies previously discussed \cite{arxiv:2401.14887v3}. Crucially, the training objective is not merely to condition the model on the passage but also to calibrate the extent to which internal, parametric knowledge can override or be corroborated by external evidence, a skill that the prompt-based source-negotiation strategies earlier attempt to elicit at inference.

While full fine-tuning can be effective for teaching these evidence-utilization policies, the immense memory and computational costs are prohibitive for many state-of-the-art LLMs. Low-Rank Adaptation (LoRA) provides a cost-effective alternative by learning a low-rank residual update for the model weights, effectively injecting or steering knowledge with a small fraction of the trainable parameters. This parameter-efficient approach has been instrumental in enabling practical RAG system deployment \cite{arxiv:2404.10981v1}. Beyond layer-level updates, the training-time design of compact plug-and-play modules can be used to connect the retriever and the LLM, such as a bridge model that is fine-tuned to optimize the connection between the components \cite{arxiv:2401.06954v2}. Similarly, a unified framework has also been leveraged to jointly fine-tune a dense retriever and a response generator with different objectives to effectively manage conversational context \cite{arxiv:2507.07030v1}. These compact module designs must confront a crucial challenge: the retriever often presents ranking and selection assumptions about ``human-friendly'' information, which is not optimally aligned with the LLM's context-processing preferences \cite{arxiv:2401.06954v2}. Such methods serve as a model-level complement to the prompt-level structural justifications for highlighting evidence, aiming to encode the desired text format and attention allocation preferences directly into the learned modules.

The field is witnessing a shift from simple likelihood maximization to explicitly teaching the model what to prefer. The core insight is that fine-tuning should increase the awareness of retrieval relevance directly in the training objective, thereby encouraging the model to produce answers that are demonstrably based on the evidence, and increasingly, to correctly decide when to ignore evidence in favor of its own parametric knowledge. Indeed, foundational work \cite{arxiv:2401.06954v2} argues that effective grounding requires a careful "chaining" of supervised and reinforcement learning to train a bridge, optimizing the connection between the retriever and the LLM through the preference space. Moreover, the performance of RAG systems is highly sensitive to the hyperparameters of both the retriever and the generator, which necessitates combined and efficient optimization of both components \cite{arxiv:2605.08333v1}. In line with this, multi-stage training paradigms have also been developed to enhance knowledge integration, for example with parametric RAG (PRAG), which incorporates knowledge via parameter updates rather than token-level attention \cite{arxiv:2501.15915v1}. Such methods help to mitigate knowledge conflicts and improve robustness to retrieval noise \cite{arxiv:2510.12668v2}. By explicitly modeling the preference between grounded and internal evidence during training, these methods aim to leave an agent less dependent on the heuristic, instruction-based prompting needed in the previous subsection when handling knowledge conflict, instead baking a stable hierarchy between sources into the model's weights.

A promising frontier is the development of dynamic knowledge injection and joint optimization systems for RAG components in an integrated manner. The entire alignment process is not about new knowledge injection but rather about optimal representation of a knowledge space where parametric and non-parametric knowledge have complementary roles. In future work, we should expect to see more sophisticated and adaptive parameter-efficient methods to dynamically modulate internal and factual knowledge based on the confidence of the retriever and the internal efficacy of the model, moving beyond the current static nature of most adapters \cite{arxiv:2510.12668v2}. This will further blur the line between efficient tuning, inference-time interventions, and prompt design. Rather than treating inference as a separate layer, the robustness and attention-preference policies forged during this training stage directly enable the lightweight, training-free decoding and attention re-weighting techniques described next; a model trained to have dense, yet selectable, evidence grounding better proxies for the mutual-information contrasts and logit adjustments measure. This training sets the fundamental capabilities to create increasingly grounded, reliable, and effective generation systems by laying the foundation for the emergent, agentic control of evidence discussed in the later prompt engineering subsection, while also laying the groundwork for the targeted interventions that will be explored next.


\subsection{Training-Free and Decoding-Time Interventions for Faithful Generation}


The training-free and decoding-time intervention paradigm represents a critical middle ground between purely parametric generation and full fine-tuning, offering a means to enhance faithfulness without the considerable computational expense of model weight updates. These methods operate exclusively at inference, manipulating the probability distribution of the decoder or its attention mechanisms to prioritize retrieved evidence. This approach is particularly appealing in large-scale deployment scenarios where a single base model is shared across many tasks and may not be fine-tuned for a specific knowledge base.

A dominant class of these interventions is \textbf{contrastive and ensemble decoding}, which seeks to recalibrate token probabilities by leveraging the discrepancy between what the model generates when conditioned on external evidence versus its own parametric priors. The theoretical underpinning draws from information theory: the mutual information between the generated sequence and the retrieved context should be maximized to ensure the output is grounded. By contrasting the output distribution of a context-augmented model \(P_\theta(y|x, c)\) with that of the model without context \(P_\theta(y|x)\), methods can re-weight the final logits to favor tokens that are statistically more probable in the presence of evidence. This approach effectively mitigates the model's over-reliance on its internal biases, a primary source of hallucination when external facts conflict with parametric memory. In scenarios involving complex, long-context reasoning, the efficacy of these contrastive methods often hinges on structured evidence presentation, as shown by techniques that organize provided documents in a logical sequence or knowledge graph to mitigate RAG failure points \cite{arxiv:2504.15909v2}. Related strategies treat the external and parametric distributions as separate streams and perform ensemble decoding, where a positive document representing the gold answer is contrasted with negative or irrelevant documents to sharpen the contextual signal \cite{arxiv:2503.04338v2}. Constructing these methods with a rigorous evaluation lens, using high-quality retrieval relevance annotations provided by human annotators, proves essential to measuring and optimizing guided reasoning quality in complex real-world scenarios \cite{arxiv:2601.08620v2}.

Constrained decoding and context-guided decoding offer another angle: rather than contrasting distributions, they alter the sequence generation by imposing hard or soft constraints derived from the retrieved evidence. Grounded score-based methods fall into this category, where the likelihood of a candidate token is adjusted based on the probability distributions computed over the retrieved context chunks. Test-time decoding strategies, such as majority voting across augmented queries, represent a simpler, training-free framework that uses consensus generation to surface robust, well-supported answers \cite{arxiv:2504.15909v2}. In addition, as context windows expand to accommodate more retrieved evidence, the risk of losing focus on critical spans increases. This motivates the integration of memory-augmented decoding strategies that selectively compress external information in the key-value cache to maintain traceability and reduce computational load within that expanded memory \cite{arxiv:2502.15734v1}. Moreover, the model can be engineered to detect subtle informational dominance between consecutive decoding memories, naturally de-emphasizing irrelevant spans without re-training, aligning with that need to manage long-context input \cite{arxiv:2410.07590v1}.

Finally, a low-cost class of training-free interventions uses plug-and-play modules attached to the decoder to modulate attention heads. These modules are optimized to "re-weight" the strength of attention heads based on the evidence's relevance, effectively denoising long contexts on-the-fly \cite{arxiv:2506.00054v1}. These interventions align to use query-oriented and knowledge-aware signals to refine retrieval and filtering in a training-free manner. This process removes the need for large gradient updates by using a secondary component that acts as a gate to filter out noise, ensuring that only the most relevant evidence influences generation \cite{arxiv:2309.15217v1}. By combining uncertainty estimation and embedding similarity to score and rank document correlation, these adaptive mechanisms function analogously to a retriever that has been augmented to be robust against noisy retrieval, offering both accuracy and stability improvements when incorporated into RAG pipelines.


\subsection{Long-Context Utilization, Context Compression, and Memory Augmentation}


The integration of retrieved evidence into generation confronts a fundamental bottleneck: the finite context window of LLMs. As retrieval scales to meet the demands of complex queries, the sheer volume of candidate documents---often dense, redundant, and heterogeneous---threatens to overwhelm the model's capacity for focused attention. This challenge is not merely a matter of input length; it is a question of \emph{effective} context utilization, where the signal-to-noise ratio degrades with increased context size, leading to the "lost-in-the-middle" phenomenon \cite{arxiv:2410.12837v1}. While simply extending context windows offers a partial solution---such as the Cache-Augmented Generation (CAG) paradigm, which preloads a static knowledge base to bypass dynamic retrieval \cite{arxiv:2412.15605v2}---this approach becomes computationally prohibitive and increasingly susceptible to attention dilution as corpora scale. More critically, the decoding-time and attention-level interventions discussed in the previous subsection treat context as a per-query constraint at the point of generation, yet they presume that the evidence itself is already presentable in a manageable form. The methods presented here complement that paradigm by addressing \emph{how} evidence can be maximally compressed, cached, or reorganized to reduce the noise and cognitive load that second-guess these generation-time interventions in the first place. Moreover, such context preparation establishes a foundational infrastructure that becomes a dynamic decision point in the adaptive, agentic architectures explored in the following section.

A first line of attack involves compressing the retrieved information \emph{before} it is presented to the generator. Token-level compaction is crucial for reducing computational overhead while preserving the critical evidential payload \cite{arxiv:2404.10981v1}. Innovations like the Tri-RAG framework address this by transforming retrieved evidence into a structured triplet format---Condition, Proof, Conclusion---which creates reasoning-aligned context that is significantly more token-efficient and semantically precise than concatenated raw text \cite{arxiv:2604.12610v1}. For systems orchestrating multiple retrieval steps, frameworks like TeaRAG compress both the \emph{retrieval content} and the \emph{reasoning steps} themselves, leveraging graph traversal to highlight key knowledge and using a token penalty within Direct Preference Optimization to reduce reasoning verbosity \cite{arxiv:2511.05385v2}. By focusing on salient information, these methods directly mitigate the risk of overwhelming the generator with irrelevant details, thereby increasing the signal-to-noise ratio that decoding-time interventions rely on to recalibrate token probabilities.

Beyond content-level compression, a second strategic line focuses on architectural and attention-level interventions to manage the massive context. The fundamental trade-off between information richness and positional attention limits is central to this issue, a problem that demands improvements in both context integration and the efficiency of long-context models \cite{arxiv:2506.00054v1}. Hardware-aware innovations like FunnelRAG employ a coarse-to-fine progressive cascade that significantly reduces time overhead without sacrificing retrieval performance \cite{arxiv:2410.10293v3}. For long conversations or extensive documentation, a persistent and dynamic system is needed to manage the information visible to the model, enabling evolving knowledge selection and redundancy removal over time \cite{arxiv:2603.01783v1}. These architectural strategies are crucial because they make the context an evolving, dynamic resource rather than a static overload, and they prefacet the dynamic retrieval-generation synergies discussed in the next section by ensuring that the pre-prepared context remains focused and fit for iterative use. The decoding process itself can also be augmented to self-correct or stop, mitigating the effects of imperfect retrieval \cite{arxiv:2505.02811v2}---preparing for the autonomous decision-making seen in dynamic RAG, where the system must decide whether its internal output satisfies the query or requires a new round of retrieval.

In sharp contrast to dynamic compression, another paradigm prioritizes \emph{efficiency} via aggressive \emph{caching} rather than compression. RAGCache and ProphetKV represent crucial advances at the system level of context management, recognizing that cache organization is a primary latency bottleneck in practice. By organizing and reusing the intermediate states of retrieved knowledge in a hierarchical memory, they can drastically reduce time-to-first-token and improve throughput \cite{arxiv:2404.12457v2}. ProphetKV pushes this further by selectively recomputing the highest-utility tokens---those aligned with the user's query---from cache to preserve accuracy without full recomputation cost \cite{arxiv:2602.02579v3}. These methods implement the memory-augmented environment and necessary practical basis that allow the later adaptive and iterative retrieval agents to operate efficiently without \emph{suffering} from the latency bottle-being that would otherwise interrupt their feedback loops.

Looking forward, the future of long-context RAG lies in orchestrating these principles---integrating semantic precision of structured retrieval with memory-driven architectures \cite{arxiv:2503.04338v2}, and injecting them into the very fabric of beyond-decoders. The shift is toward designing \emph{selective} relevance engines rather than naive truncation, and evolving from a static index to a a "stateful" RAG system that continually learns from its own queries and feedback. This is precisely what enables the dynamic agents in the following subsection: while their rich reasoning loops are only viable if the core context is efficiently managed and cached. As LLMs become increasingly capable of extended reasoning via multi-agent and CoT frameworks \cite{arxiv:2504.15909v2,arxiv:2507.09477v2}, context management** must evolve to serve as a dynamic, selectively refreshed and attention-aware resource---transforming the limitation of a constant window into a continuously moderated, adaptive infrastructure fully grounded in both attention and memory.


\subsection{Dynamic and Adaptive Retrieval with Generation-Aware Mechanisms}


The conventional view of retrieval-augmented generation as a static, pipeline-style process---where a query is transformed, a fixed number of documents are fetched, and the generator consumes the amalgamated context in a single pass---is being fundamentally challenged by a new class of methods that treat retrieval as a dynamic, generation-aware process. This paradigm shift, central to what is increasingly termed "Dynamic RAG" \cite{arxiv:2506.06704v1}, posits that the retrieval decision should not be a pre-hoc heuristic but an integral and interactive component of the generation itself. In these systems, the LLM's internal state---ranging from activation patterns to token-level entropy---serves as a continuous, real-time gauge of information sufficiency, enabling the model to autonomously determine \emph{when} to retrieve, \emph{what} to retrieve, and, critically, when to stop generating. This approach optimizes both answer accuracy and computational efficiency by removing the reliance on external pipeline orchestrators or rigid, pre-defined rules.

A foundational principle of this dynamic paradigm is the interrogation of the model's latent representations to determine retrieval necessity. Instead of indiscriminate retrieval, models can use "introspection layers" or lightweight probers that analyze hidden states to identify a "failure state" or knowledge gap. For instance, Skill-RAG couples a hidden-state prober with a prompt-based skill router; upon detecting a failure state, it activates one of several retrieval skills---such as query rewriting or question decomposition---to correct the query-evidence alignment before/after a subsequent generation attempt \cite{arxiv:2604.15771v3}. Research probing LLM representations has demonstrated that their hidden states encode not just the content of the output but also a verifiable signal of the source of that knowledge (internal \emph{vs.} external), a signal that can be employed for "knowledge check" and filtering of noisy context \cite{arxiv:2411.14572v1}. Similarly, gating strategies and confidence-based mechanisms trigger retrieval based on an analysis of user intent and ongoing contextual needs and identification of "irreducible cases," as formalized in frameworks like BRIDGE \cite{arxiv:2505.17118v2}. However, the most intelligent systems do not \emph{solely} rely on pre-generation probes, but rather interleave retrieval with the generation process itself.

This feedback-driven view has led to the development of iterative retrieval-generation synergies, where the partially generated output from the LLM is leveraged as a rich, context-specific query for the initial/next retrieval step. A robust school of thought follows the philosophy of "improving retrieval with generation," suggesting strong performance can be achieved by iteratively using a model's initial output to retrieve more relevant knowledge, which in turn helps generate a better output, a method exemplified by Iter-RetGen \cite{arxiv:2305.15294v2}. These methodologies, often characterized as rolling retrieval, create a "retrieval-generation synergy," in which the model output shows what might be needed to finish a task. Building upon this, Self-RAG trains a single arbitrary LM to emit special "reflection tokens" that allow it to \emph{adaptively} decide on-demand whether to retrieve, and to critique its own generations for usefulness and factuality against the retrieval passages, thereby allowing the system to "tailor its behavior to diverse task requirements" at inference time \cite{arxiv:2310.11511v1}. As an alternative to this optimization, implicit search \emph{agents} have been trained using reinforcement learning with complex rewards. For example, R1-Searcher++ introduces a two-stage process of SFT cold-start followed by reinforcement learning for "Dynamic Knowledge Acquisition," in which its reward mechanism distinctly incentivizes the use of both internal knowledge and external search engines, allowing for the "efficient orchestration of retrieval" \cite{arxiv:2505.17005v1}.

The most advanced form of this dynamic behavior integrates explicit agentic planning into the retrieval sequence. Autonomy manifests in the ability to plan with reasoning, not just in terms of \emph{when} to search, but also in planning the logic of the search, often formalizing a Markov Decision Process (MDP) to facilitate complex decision-making. DeepRAG, for instance, models retrieval-augmented reasoning as a Markov Decision Process, iteratively decomposing queries to dynamically decide whether the model should retrieve external knowledge or rely on parametric reasoning at each step \cite{arxiv:2502.01142v2}. However, these iterative steps can lead to multi-token reasoning loops that bloat cost; to counteract this, TeaRAG uses a multi-faceted approach that combines the token efficiency of "knowledge association graphs" with 'Personalized PageRank', and an iterative, process-aware optimization method (IP-DPO) that learns to balance efficiency and accuracy \cite{arxiv:2511.05385v2}. In this vein, *Auto-RAG emphasizes the model's ability to "autonomously adjust the number of iterations based on the usefulness of the retrieved knowledge" via multi-turn dialogue interplay with the retriever \cite{arxiv:2411.19443v1}.

In conclusion, the future of RAG is not in designing better static pipelines, but in creating synthetic models that are more adaptive and self-aware. The current trajectory indicates a synthesis between high-level, search-driven exploitation of parametric and non-parametric memory. These mechanisms replace one-shot "frozen parameters" with a robust and dynamic system that internalizes knowledge checks at various levels, deciding what to trust and when to rely on latent memory \cite{arxiv:2410.05162v1}.


\section{Training, Alignment, and Optimization of Full Systems}



\subsection{Joint End-to-End Optimization of Retriever-Generator Pipelines}


The joint optimization of retriever and generator components represents one of the most intellectually demanding yet practically consequential challenges in the RAG paradigm. Early RAG formulations, such as the foundational sequence-to-sequence framework with a differentiable dense retriever, established that end-to-end training is feasible when the retrieval process is treated as a latent variable \cite{arxiv:2005.11401v4}. However, the inherently discrete nature of document selection creates a non-differentiable bottleneck between the retriever's continuous scoring function and the generator's token-level objective, necessitating sophisticated gradient estimation techniques.

Policy-gradient methods have emerged as the dominant framework for bridging this divide. By treating the retriever as a stochastic policy that selects a document distribution conditioned on the query, the generator's likelihood can be employed as a reward signal, enabling REINFORCE-style updates. This approach is exemplified in training paradigms where the generator, often initialized from a pre-trained checkpoint, provides scalar feedback that propagates through the sampled retrieval path \cite{arxiv:2005.11401v4}. The primary advantage of such reinforcement learning formulations lies in their model-agnostic nature: the retriever can be optimized against any downstream generation objective without requiring the generator to be differentiable with respect to retrieval parameters. For instance, concurrent work has shown that a Markov Decision Process formulation of retrieval-augmented reasoning---where each decomposition step decides whether to retrieve or generate---can be effectively optimized to improve both retrieval efficiency and the accuracy of the final answer \cite{arxiv:2502.01142v2}. Furthermore, beyond policy-gradient, expectation-maximization-like algorithms provide an alternative optimization trajectory; treating the retrieved documents as a latent variable and performing iterative refinement enables a clean decomposition of the learning problem into explicit retrieval and generation steps \cite{arxiv:2005.11401v4}. This result suggests that by marginalizing, rather than merely sampling, more robust parameter updates can be obtained.

A key challenge in the policy-gradient framework is the high variance inherent in the credit assignment from the generator's scalar reward back to the discrete selection. Standard solutions include subtracting a baseline reward (e.g., the reward from an average document or from the parametric-only generation) to reduce unnecessary variance. To further stabilize training, many modern systems employ a two-stage learning scheme: a pre-training of the retriever using well-defined supervised contrastive objectives, followed by a constrained fine-tuning phase where the retriever's weights are updated against the downstream generator's score under a straighter, path-wise objective \cite{arxiv:2005.11401v4}. This hybrid approach is precisely where the core tension in joint optimization lies: a fully actor-critic setup is computationally intractable in huge search spaces, but a purely pre-trained retriever remains misaligned with the generator's semantics. The introduction of a plug-and-play search loss such as Minimum Spanning Tree, and its many variants (e.g., Gumbel-Top-k and straight-through estimators), directly address the obstacle of discrete back-propagation. Also, techniques that treat retrieval as a learned function not from \emph{clicks} but from generative feedback offer a promising axis. This is conceptually proven by adversarially or heuristically fine-tuning a generator so it can be used as a frozen evaluator of the retriever: even in a completely end-to-end training setup, "generator-as-teacher" signals can distill, from the coordination friction. An interesting yet newly suggested line of research is to incorporate retrieval robustness as a training objective, where both retriever and generator components' failures are jointly considered. In this route, the loss is not only the amount of desired document likelihood, but also a robustness term that penalizes downstream generation under imperfect retrieval. This has grown hand in hand with the insight that in the pipeline, both the retriever and generator are prone to error propagation. By training a system where components can directly overrule or communicate ignorance about the quality of the retrieved document, we overcome the assumption that the model must operate without mutual awareness, which greatly improves robustness to conflicting evidence \cite{arxiv:2410.07176v2}. This solution always aligns the two components under a shared objective. Future directions will likely see the formulation of more principled objectives that directly optimize the final evaluation, alongside a dynamic learning of the distribution of policies on a Markov Decision Process \cite{arxiv:2502.01142v2}.


\subsection{Alignment, Instruction Tuning, and Multi-Task Learning}


The alignment of retrieval-augmented generation (RAG) systems extends beyond the joint optimization of individual components discussed previously; it necessitates supervising the integrated behavior of the retriever-generator dyad. While joint optimization focuses on aligning the retriever's selection with the generator's objective through gradient-based coordination, this subsection explores the critical paradigms of instruction tuning and preference optimization, which have become indispensable for instilling desired behaviors into the combined system---primarily, the capacity to faithfully utilize external evidence while remaining robust to the inevitable noise and conflicts present in retrieved corpora \cite{arxiv:2312.10997v5}. This training regime directly confronts the retriever-generator semantic gap \cite{arxiv:2507.07030v1}, ensuring the generation component can reliably interpret and act upon the context provided by its retrieval counterpart.

A foundational approach involves supervised instruction tuning on curated (query, context, answer) triples. This approach teaches the generator to condition its outputs on retrieved evidence, effectively bridging the gap between internal parametric knowledge and non-parametric inputs \cite{arxiv:2405.06211v3}. However, within a RAG context, the generator must not simply use the context; it must learn how to handle the spectrum of retrieval quality---from highly relevant and clean evidence to irrelevant, noisy, or even contradictory passages. To this end, recent research highlights the value of employing the instructor's objective to establish large-scale multi-task datasets that expose the model to various context conditions as a primary objective of the tuning process \cite{arxiv:2401.06954v2}. A crucial innovation in this space is the design of training signals that teach the model to be robust to sub-optimal retrieval, thereby preventing the retriever from becoming a single point of failure in the pipeline \cite{arxiv:2506.00054v1}.

However, simply training on (query, context, answer) triples does not explicitly guide a model to navigate the distinction between internally memorized information and externally supplied evidence---especially where this distinction is critical. Therefore, alignment approaches based on preference optimization have emerged as a powerful mechanism to shape system behavior \cite{arxiv:2411.19443v1}. Methods like Direct Preference Optimization (DPO) have been adapted for RAG, where the system is trained on a preference dataset whose contrasts desired and undesired outputs, enabling the model to internalize the nuances of grounding. For instance, rather than merely learning to "use" the context, preference optimization methods can be designed to penalize hallucinated content using retrieved evidence, thereby aligning the system with objectives of increased faithfulness and improved citation integrity \cite{arxiv:2401.06954v2}. This process is particularly important for handling knowledge conflicts, where retrieved non-parametric knowledge contradicts the model's stored parametric memory. The goal of alignment, in this case, is to teach the model to prefer correct grounding, reinforcing the insight that in conflict scenarios, a RAG model should undertake the contextual evidence more heavily than its parametric memory \cite{arxiv:2410.05162v1}.

Preference optimization becomes a valuable tool for bridging the fundamental reward gap between retrieval and generation. Furthermore, the training objective can be tuned not only to improve answer quality but also to cultivate robustness under persuasion. By generating "noisy" or "unhelpful" contexts as negative samples, the aligned generator can learn to either selectively ignore sub-optimal evidence or explicitly override it, extending the scope of alignment into granular, task-level generalization across a wide range of retrieval failures \cite{arxiv:2404.10981v1}.

The frontier of this area concerns how the system learns to reconcile retrieval limitations with appropriate language generation while gracefully integrating the preceding training signals. The emergent trend is to move away from static, one-shot supervision towards iterative, honest feedback loops where the training set is driven by retrieval quality, challenging the system to both use and override retrieved evidence based on a principled curriculum---one that equally trades off parametric and non-parametric memory \cite{arxiv:2005.11401v4}. This cross-cutting approach converts what is essentially a data challenge into an effective alignment strategy that complements the system-level coordination objectives. Importantly, this synergy directly anticipates the following subsection on distillation: preference-aligned teachers, trained with the ability to selectively ground evidence, serve as far more effective and reliable data generators for downstream students. Instead of propagating flawed synthetic signals, these aligned teachers can distill their grounding quality into compact student models, enabling cost-efficient deployment without heuristics propagation. Ultimately, this trajectory---from joint optimization, to system-level preference alignment, and finally to distillation---points towards a cohesive framework for building RAG systems that are both robust and practical, where each stage reinforces the safe and grounded behavior of the final integrated system.


\subsection{Distillation and Synthetic Data for System Training}


The prohibitive cost of jointly optimizing large-scale retrievers and generators has catalyzed the emergence of distillation and synthetic data generation as cornerstone techniques for training RAG systems. Rather than relying on expensive human annotations or full-model gradient updates, these approaches leverage large, often cumbersome teacher models to produce supervisory signals that train efficient components, addressing several coordinating challenges including data scarcity, domain adaptation, and computational overhead \cite{arxiv:2405.07437v2,arxiv:2404.10981v1}. As the field has matured, a methodological divide has emerged: one focused on distilling knowledge into the retriever, and another concerned with leveraging synthetic data to enhance the generator's grounded behavior.

A primary objective of distillation in RAG is to train compact student retrievers capable of approaching the performance of their larger counterparts. Traditional dense retrievers often struggle to capture the complex semantics required for ranking in knowledge-intensive tasks, leaving a performance gap in comparison to LLM-based or cross-encoder rankers. Distillation methods address this by using the LLM's per-passage output probabilities or its latent representations as a richer supervisory signal compared to binary labels. Specifically, the teacher's conditional likelihood over a set of candidate documents can be treated as a soft target, enabling the student to learn the relative importance of each document in a list-wise fashion---a paradigm that produces more robust and fine-grained retrieval scores than point-wise discriminant training alone. Consequently, this two-stage paradigm---pre-training a general student, then distilling with an LLM oracle---produces lightweight retrievers with a performance level previously unattainable at their scale, significantly reducing the inference costs and latency associated with joint systems \cite{arxiv:2405.07437v2,arxiv:2507.18910v1,arxiv:2503.14649v2}.

Parallel to retrieval optimization, there is a growing class of techniques that utilize synthetic data to instruct and align the generation module. Considering the extensive cost of manually curating (query, document, response) triples, particularly for low-resource domains, the field has turned toward prompting large teacher LLMs to produce training data \cite{arxiv:2312.10997v5,arxiv:2504.08893v1,arxiv:2506.00054v1}. These teachers can generate ground-truth answers, auxiliary chain-of-thought reasoning sequences, and even diverse and difficult query variants. This corpus is then used for supervised instruction fine-tuning of smaller student generators, allowing a system to be adapted to domain-specific vocabulary and heterogeneous data structures without the reliance on costly human annotations. This process also creates a cheaper-as-you-go cycle to generate data and provide valuable support for domain adaptation and private deployment.

A comparative look at these strategies reveals a critical insight; data synthesis can be used as a source for direct imitation. Two dominant paradigms have emerged in this area: distillation of the retrieval-aught teacher's workflow and the use of distillation for consistency. The former is an efficient means to remove the index runtime requirement at deployment; in this case, a cumbersome RAG pipeline generates a set of teacher outputs from queries, and the student model is trained to predict those outputs solely from the query, internalizing the reasoning and grounding the model learned from the retrieved passages. The other, perhaps more subtle, path is distillation for internal consistency, which aims to make the final generator robust to variations in the retrieved context---for instance, by teaching a model to output consistent predictions regardless of the order of the retrieved chunks, thereby reducing 'lost-in-the-middle' errors. This rank-distillation set of techniques significantly reduces the negative effects of noisy retrieval \cite{arxiv:2501.07391v1}.

Emerging trends point to a convergence of these avenues. The next generation of systems will likely blur the boundaries between teacher and student through knowledge-sharing loops and agentic data generation techniques, where the LLM uses tools to create data fromand in concert withthe retrieval index and its own feedback \cite{arxiv:2412.05838v1,arxiv:2504.15909v2}. The primary challenge in the production of synthetic external knowledge lies in mitigating hallucination propagation, where teachers may generate plausible but incorrect data, which then gets injected into downstream tasks. To ensure these distillation approaches gain broad production acceptance, the mitigation of teacher-hallucination in the generation stage will require robust validation and data augmentation. The principal focus will be on using them to improve search effectiveness and accuracy, while general-purpose applications are pushed toward ever-lower costs without quality degradation, a crucial step for practical implementations \cite{arxiv:2507.18910v1,arxiv:2502.15734v1}.


\subsection{Continual and Dynamic Learning for Evolving Knowledge Environments}


The operational lifespan of retrieval-augmented generation systems is fundamentally constrained by the assumption of a static knowledge corpus. In practice, real-world deployments face non-stationary environments where information is continuously created, updated, and rendered obsolete, necessitating a shift away from the costly paradigm of full re-training and complete re-indexing. This subsection addresses the critical challenge of maintaining system relevance over time, exploring incremental indexing, adaptive memory management, and continual learning from feedback---the architectural and algorithmic innovations that enable RAG systems to co-evolve with the landscapes they are intended to model. While the previous subsection examined how distillation and synthetic data can train efficient components from static supervision, the techniques discussed here confront the complementary problem of \emph{ongoing} adaptation: how a system, once trained, can remain accurate and useful as both the knowledge base and the query distribution evolve over time, without revisiting the expensive, large-scale optimization procedures previously described.

The most immediate challenge in dynamic environments is maintaining an up-to-date knowledge store without the prohibitive cost of reconstructing index structures. Traditional vector search indexes are largely batch-oriented; inserting or updating a document typically triggers recomputation over a broad index neighborhood. To circumvent this, incremental indexing strategies decompose the dynamic update problem into coarse-grained management and fine-grained integration. A notable approach, GAM-RAG, introduces a training-free framework that accumulates retrieval experience from recurring queries within a lightweight hierarchical index \cite{arxiv:2603.01783v1}. This framework leverages sentence-level feedback from successful retrieval episodes to update a semantic co-occurrence graph, where links represent potential usefulness for a class of queries rather than static relational semantics. The critical design is a Kalman-inspired gain rule that jointly updates memory states and uncertainty estimates, applying fast updates to novel signals and conservative refinements to stable or noisy memories, implicitly serving as a gated memory channel. Such an approach directly addresses a fundamental limitation of traditional RAG: a static index remains blind to new evidence but also fails to account for how queries continually shape which portions of the knowledge space are relevant. This mechanism illustrates a broader theme of this subsection: moving from offline, one-time distillation of knowledge into a model or index, toward online, episode-driven updates that are perpetually refined through interaction. A related class of approaches focuses on schema-based learning and hierarchical cognitive models, with a broader aim of bridging single-query semantic fetches to multi-query long-horizon contexts \cite{arxiv:2506.10408v1}. Rather than ending with a single retrieval index, these designs update the representation with every successful interaction, making retrieval "smarter" with every episode.

Building on the update of the index structure, dynamic adaptation must also address the distribution shift of the queries themselves and the adaptability of the model components. While the index is the primary repository of knowledge, systems remain relevant only if both retrieval and generation subcomponents can leverage novel data patterns. Some frameworks treat retriever adaptation as a form of sequential decision-making, using multi-armed bandit selection to dynamically choose retrieval strategies in response to evolving query complexity \cite{arxiv:2501.00332v1} \cite{arxiv:2412.01572v4}. This perspective positions retrieval not as a static algorithm but as an online learning problem where the system explores different retrieval strategies to balance accuracy and computational budget in real time. On the generation side, continual adaptation demands a careful balance with catastrophic forgetting in downstream supervised fine-tuning \cite{arxiv:2509.03934v1}. The simultaneous updating of a retriever and generator is increasingly orchestrated through parameter-efficient fine-tuning modules (which the following sections will elaborate) or preference optimization that penalizes reliance on stale evidence. However, one must resist the notion that dynamic updates should be global to the LLM itself---a key system design choice is to keep the internet of the LLM intact while providing a tool chain to interact with the dynamic index. In these cases, the orchestration layer or agentic element often hosts the short-term memory necessary for multi-step or multi-turn exploration \cite{arxiv:2510.13910v2,arxiv:2412.05838v1}, emphasizing dynamic mechanism for controlling the point of external knowledge entry. This design philosophy directly corresponds to the preservation of the foundational model's generality, a theme that the following subsection will address in depth: just as distillation can make aligned behavior computationally feasible, here we see that leaving the weights unchanged and instead adapting the surrounding data infrastructure serves as an equally crucial strategy for maintaining model utility in shifting environments.

Regardless of the implementation details, modern dynamic RAG systems must balance three constraints: an ever-growing knowledge space, the finite context or compute budget of the generation step, and the requirement to maintain temporal coherence. Continual learning from feedback loops, combined with a memory that consists of summarized, high-level abstractions of retrieved episodes, are emerging as paramount. Looking forward, the canonical future of retrieval-augmented generation lies in deploying systems that can \textbf{proactively} manage their own memory---deciding not only what to retrieve in the short term, but what to remember in the long term. This entails a shift from static indexes toward \textbf{indexes as evolving agents}, that remain relevant to the task at hand and adapt comfortably within a multi-query problem-solving workflow. As the following subsection will show, these evolving indexes must be complemented by operational layers that customize how the model consumes them---without sacrificing the general-purpose capabilities of the LLM. The challenge ahead lies in deriving robust theoretical frameworks that reason about the stability-plasticity trade-off to yield definitive inference-time bounds, while effectively and safely pruning obsolete memories without sacrificing historical context. Hence, the field is increasingly heading toward the "index as an evolving agent", not merely an auxiliary database supported over ever-changing data domains.


\subsection{Adaptation and Parameter-Efficient System Optimization}


The deployment of Retrieval-Augmented Generation (RAG) systems in specialized domains frequently encounters a fundamental tension: the need for task-specific customization against the prohibitive cost of full-model fine-tuning. Full fine-tuning of large language models (LLMs) is not only computationally expensive but also risks catastrophic forgetting of general capabilities. This subsection examines parameter-efficient and adaptation strategies that preserve the foundational model's weights while customizing system behavior for domain-specific or task-specific requirements, covering lightweight adapters, low-rank matrices, learnable prompts, and test-time adaptation methods.

A primary line of work leverages parameter-efficient fine-tuning (PEFT) techniques to insert lightweight, trainable modules into the frozen retriever and generator components. Low-Rank Adaptation (LoRA) and its variants have become the de facto standard for adapting LLMs to new domains, operating by learning low-rank decomposition matrices that are added to the model's existing weight matrices. In the RAG context, \cite{arxiv:2501.18365v1} extends this idea to explicitly address realistic deployment challenges by proposing robust Fine-tuning tasks that enhance LLM resilience against retrieval defects---including noisy, irrelevant, or counterfactual information---demonstrating that PEFT can target downstream system behaviors beyond standard task performance. The Astute RAG framework \cite{arxiv:2410.07176v2} further illustrates the necessity of adaptation at the component level, for understanding the inherent tension between internal and external knowledge that becomes a bottleneck for robust feedback retrieval. Beyond adjusting the generator, adaptation strategies also extend to the fusion interface. Plug-and-play modules that sit between the retriever and generator can be trained in multi-stage paradigms to extract, compress, or reorganize retrieved chunks into query-specific formats preferred by the frozen LLM. These modules function as deterministic intermediate representations, parameterizing the interface between modules without altering the base model weights.

Prompt-based and in-context adaptation offer an even lighter-weight alternative to PEFT, particularly for scenarios with minimal data or a need for rapid deployment cycles. By optimizing a set of soft prompts or hard examples on top of frozen model weights, a RAG system can be steered towards a new domain while also being taught how to more effectively utilize the provided context, including learning specific formats for evidence handling. However, a critical limitation of some benchmarked methods is their sensitivity to data distribution and retrieval settings, which PEFT methods often mitigate more stably. The broader design space also includes training-time intervention, where the reader adapts to fit a specific workflow---for instance, using alignment techniques to enhance the interplay between the model's parametric knowledge and retrieved evidence. Beyond fundamental fine-tuning, exceptional efforts, such as the TCR framework \cite{arxiv:2601.06842v1}, explore how lightweight, soft-prompt-style methods can be combined with weights (e.g., SNR-based weighting) to handle specific, noisy external sources. Furthermore, adaptation can be applied to the retrieval strategy as well: frameworks like MBA-RAG \cite{arxiv:2412.01572v4} propose a reinforcement learning-based bandit approach that dynamically selects a retrieval strategy based on query complexity, adaptively balancing accuracy and efficiency without modifying the underlying model weights, highlighting how adaptation to query- or domain-specific demands is an operational task as much as a weight-tuning one. Critically, the choice of PEFT over full fine-tuning also serves as a question of trade-off: frameworks such as \cite{arxiv:2504.15909v2} emphasize the benefits of task-specific reasoning against the generality of the foundational model. Test-time strategies further resolve this via related methods, allowing for per-instance adaptation to unique data distributions. For instance, in federated search settings, lightweight routing can adapt a RAG system on-the-fly to select relevant sources, thereby improving precision and reducing the risk of retrieval errors \cite{arxiv:2502.19280v2}.

In synthesis, the field is moving from monolithic retraining toward modular, state-aware, and dynamic adaptation strategies. The growing reliance on PEFT and test-time methods is being driven by a critical need for resource efficiency and the desire to preserve the generality of the foundational model. Future directions for this area include the development of generalized, composable adapters that can be switched on the fly to serve multiple sub-competencies within a single system and the use of adaptation signals derived directly from language model internal states to fine-tune overall robustness. Achieving the best of both worlds---highly specialized performance with world-level generality---will require not just innovation in the structure of the adapters, but in the training objectives and in nuanced procedures through which a generic model is psychologically prepared to accept context.


\section{Evaluation, Benchmarking, and Trustworthiness}



\subsection{Component-Level and End-to-End Evaluation Metrics}


Evaluating retrieval-augmented generation (RAG) systems demands a departure from the isolated assessment of their constituent parts, as the complex interplay between retrieval and generation fundamentally shapes end-to-end performance. The axiom that a system's quality is bounded by its weakest component is particularly salient here: perfect retrieval cannot compensate for a generator that fails to integrate evidence, just as a robust generator cannot overcome an ineffective retriever. A foundational evaluation framework must therefore provide a disaggregated view of each component alongside a holistic assessment of their interaction. This distinction underscores a critical challenge in the field---the inherent tension between component-level optimization and system-level behavior. For instance, a retriever optimized solely for ranked accuracy may not align with the generator's needs, particularly in long-form generation tasks where the identification of dispersed key points is more critical than the mere provision of a ranked list \cite{arxiv:2506.20051v2}. Conversely, evaluating only generation quality obscures whether errors stem from a failure to retrieve relevant evidence or from a failure to condition on it.

The evaluation of the retrieval component has historically been grounded in traditional information retrieval metrics. Relevance-based measures such as precision@\(k\), recall@\(k\), and their rank-aware counterparts like nDCG and MRR quantify the retriever's ability to surface pertinent information. However, the utility of these metrics in predicting final RAG quality is limited. The eRAG paradigm provides a compelling insight by positing that the true quality of a retrieved document should be defined by its individual contribution to the accuracy of the final generation \cite{arxiv:2005.11401v4}. In this framework, a document is considered relevant only if its inclusion in the context leads to a more correct answer. This "downstream-aware" evaluation reveals a low correlation between standard component-level retrieval scores and end-to-end RAG performance, suggesting that future retrieval metrics should be grounded in, or calibrated against, the generator's behavior \cite{arxiv:2410.12837v1}. This shift is critical for guiding system design: it mitigates the risk of optimizing the retriever for an objective that does not translate to improved system utility, a phenomenon particularly pronounced in complex multi-hop reasoning and domain-specific tasks \cite{arxiv:2309.01431v2}.

For the generation component, a wide array of metrics exists, spanning lexical overlap, semantic similarity, and more recently, task-specific and faithfulness-oriented criteria. Traditional metrics like BLEU, ROUGE, or BERTScore assess overlap with reference answers but often fail to capture the required "attribution" quality inherent in RAG systems. Similarly, a sole focus on semantic equivalence can neglect the fact that an output might be factually accurate without being "grounded" in the provided, retrieved context. This metric deficiency has driven the shift toward task-aware evaluation, especially for long-form knowledge generation. For instance, the Long$^{2}$RAG model introduces Key Point Recall, a metric that assesses the generator's ability to extract and synthesize key information expected from the retrieved document corpus \cite{arxiv:2312.10997v5}. This metric focuses on the selection and synthesis of information rather than on fluency or language quality, as the semantic content of the generated output is the primary target \cite{arxiv:2506.20051v2}. The choice between these types of metrics, however, introduces inherent trade-offs: task-specific metrics require a reliable ground truth and often rely on a rigorous human or model-as-a-judge, whereas more universal, lexical metrics may be less costly but fail to provide the same diagnostic signal \cite{arxiv:2405.07437v2}.

The symbiotic nature of the retriever-generator interaction has given rise to integrated evaluation frameworks that treat them as a single entity. Such frameworks synthesize component-level metrics into a single ranking score, positing that the components themselves are not evaluated in isolation but for their contribution to the integrated system \cite{arxiv:2501.13958v3}. However, these component-level evaluations, combined with end-to-end metrics, can be useful for thread analysis. This is essential for system diagnosis, bridging the gap between an ideal system and its actual behavior. Furthermore, the diagnosis extends to a conceptual categorization of RAG outcomes, such as the source of a deficiency, via a confusion matrix. This categorization, mandated by the ongoing need for a structured evaluation, is an attempt to disentangle attribution, fluency, and faithfulness trade-offs in a human-centric manner \cite{arxiv:2405.07437v2}. It is also clear that there is a difference between the information required by the retriever and the information needed to prompt the generator. To provide a truly effective evaluation, the community must move towards metrics that provide fine-grained, diagnostic signals over the components, while still recognizing the inherent dependency between them and avoiding to treat them as fully independent entities.


\subsection{Benchmarking Datasets and Evaluation Frameworks}


The evaluation of retrieval-augmented generation systems has evolved from a reliance on generic natural language benchmarks toward a more nuanced ecosystem of specialized datasets and diagnostic frameworks. This shift reflects a growing recognition that static, single-hop benchmarks fail to capture the interaction between retrieval fidelity and generative reasoning under realistic conditions. While the preceding discussion established the theoretical and metric foundations for component-level and system-level diagnosis, the practical realization of these diagnostic ideals depends critically on the availability of diverse, dynamic, and realistic datasets. Modern benchmarking efforts differ along several critical dimensions: task coverage, whether they target component-level or system-level behavior, whether they support dynamic knowledge situations, and the fidelity of the retrieval pool to real-world corpus challenges.

Early evaluations drew primarily on knowledge-intensive NLP tasks such as open-domain question answering, fact verification, and language generation, establishing foundational practices for coupling retrieval to generation \cite{arxiv:2005.11401v4}. However, as RAG architectures have evolved, these benchmarks have become insufficient for isolating failure modes, prompting a movement toward intent-driven and multidimensional evaluation. To precisely diagnose each component, dedicated benchmarks now target both distinct functional capabilities and realistic system-level demands. This requires careful categorization of benchmark tasks by scope and design. For instance, emerging suites of advanced RAG frameworks are modularized to evaluate each phase---pre-retrieval, retrieval, post-retrieval, and generation---allowing direct identification of where pipeline failures occur \cite{arxiv:2412.15529v4}.

From a category analysis, standardized datacentric benchmarks are being augmented by cross-model heterogeneous sets. The research community has recognized that evaluating RAG systems solely on internal and often memorized corpora leads to overscore and memorization. This shift toward dynamic corpora is highlighted by the need to track real-world events, such as the application of the Need for Dynamic Knowledge Bench, which provides automatic, on-demand benchmarks to cut leakage and data overlap with pre-training corpora. Such controlled test frameworks ensure datasets remain updated, which is essential for accurately diagnosing retrieval accuracy and temporal freshness. Similarly, specialized challenges in dynamic test schedules are supported by frameworks like CoFE-RAG, which facilitate evaluation across a full-chain pipeline while enriching data diversity in document formats and query types, thus catching defects in chunking, retrieval, reranking, and generation \cite{arxiv:2410.12248v1}.

Another major frontier in benchmarking is the integration of generated evaluation protocols that act as arena-based interaction tests, surpassing prior static corpora. The introduction of MIRAGE (Multilingual Information Retrieval Augmented Generation Evaluation) offers a vast, diverse, real-world arena, where surrogate models are used to judge answer quality as correct (See Multi-Aspect Approach). This approach moves beyond simple golden-answer comparison and better measures trustworthiness under combinatorial events. In low-resource and multilingual settings, such automated judges offer a necessary scalability that would otherwise demand prohibitive amounts of human annotation \cite{arxiv:2405.06211v3}. As model capability advances, all-in-one holistic evaluators that integrate metrics through a recomposition strategy spanning component and system level control are possible to obtain a trustworthy, unified score \cite{arxiv:2506.00054v1}.

Yet, a critical gap persists: current evaluation frameworks may be too disconnected from real deployment constraints, especially regarding domain-specific noise, context comprehension, and runtime overhead. Though benchmarks specifically designed for domain-specific contexts such as health and finance highlight challenges not present in academic QA \cite{arxiv:2411.11895v1}, in production, metrics must also fit service-level objectives such as latency and concurrent capacity, requiring benchmarking to investigate beyond natural language to also assess the complexities and system trade-offs \cite{arxiv:2412.11854v1}. Future innovative benchmarks must align multi-facet robustness of components with highly concurrent and multi-turn interactions and deployed system constraints, and expand the safety and security threads; these datasets will be crucial in determining if the next-generation approaches for system-level accuracy under constraints are truly trustworthy. This pragmatic turn---anchoring benchmark design in real-world constraints and failure modes---directly connects to the adversarial and perturbed scenarios that the following section addresses, ensuring that evaluation not only measures well but also validates resilience against the threats that RAG systems will inevitably face.


\subsection{Faithfulness, Attribution, and Query-Level Robustness}


The evaluation of faithfulness, attribution, and query-level robustness constitutes the epistemic core of trustworthiness assessment in RAG systems, addressing the fundamental question of whether generated content is verifiably grounded in retrieved evidence rather than in the model's parametric priors. Unlike component-level metrics focused on retrieval accuracy or fluency, this dimension interrogates the truthfulness, provenance, and resilience of RAG outputs---critical properties for deployment in high-stakes domains such as medicine, law, and finance \cite{arxiv:2404.10981v1}.

Faithfulness evaluation centers on the functional consistency between a generated statement and its supporting evidence. Recent work has moved beyond lexical-similarity metrics (e.g., ROUGE, BLEU) toward entailment- and contradiction-aware frameworks that explicitly measure evidence coverage and logical alignment. For instance, the RAG Assessment framework (RAGAS) proposed reference-free assessment of answer faithfulness by decomposing generated statements into atomic claims and independently verifying each claim against the retrieved context \cite{arxiv:2309.15217v1}. This claim-checking procedure formalizes a multi-stage entailment inference: for a generated answer \( \mathcal{A} \) and retrieved evidence set \( \mathcal{D} \), the faithfulness score is computed as the fraction of statements \( s \in \mathcal{A} \) that are entailed by at least one \( d \in \mathcal{D} \): \( F(\mathcal{A}, \mathcal{D}) = \frac{1}{|\mathcal{A}|} \sum_{s \in \mathcal{A}} \mathbb{1}[96] \). Related work has further developed the RAG confusion matrix, which provides a taxonomy of RAG behaviors (e.g., correctly grounded, partially grounded, conflicting, and ungrounded) to facilitate fine-grained diagnostic analysis---a decomposition that is essential for auditing system failure modes beyond average accuracy scores.

Attribution evaluation extends faithfulness by assessing the model's ability to provide source-grounded, verifiable references---namely, whether each generated claim is accompanied by a precise, correct citation to the retrieved document. This is conceptually distinct from faithfulness, since a system may be faithful yet fail to traceably attribute its content, or accurately cite irrelevant documents. Recent frameworks frame attribution in terms of semantic alignment between a generated claim, its surrounding text, and the cited document snippet. Methods such as the Answer Attribution Engine use an LLM-as-a-judge prompt to assess whether a claim can be logically derived from an associated citation, yielding high correlation with human annotators \cite{arxiv:2312.10997v5}. Crucial advances have also addressed the problem of \emph{citation bias}, where RAG systems favor generic or positionally biased sources rather than the most informative ones, and in this context propose using spectral or attention-based probes to detect and re-rank source preferences \cite{arxiv:2506.00054v1}. Hierarchy attribution frameworks further distinguish direct, synthesized, and derived citation forms, and relate these to the reliability of user trust \cite{arxiv:2503.18016v1}.

A crisis emerging in attribution evaluation is the need for \emph{non-generative side measures} that do not rely on second LLM judges---an approach prone to propagating judge-side hallucinations and compounding uncertainty. This has led to information-theoretic scoring, such as measuring mutual information between the output text and the cited evidence, and verification via \emph{counterfactual document substitution}: replacing the cited passage with a distractor of comparable surface similarity and measuring prediction degradation. The mitigation pipeline thus includes both prompt-based judge pipelines and representation-based classifiers that probe a generation model's internal hidden states to predict per-claim faithfulness, without requiring an auxiliary generation \cite{arxiv:2411.14572v1}.

Query-level robustness constitutes an orthogonal but major axis---concerned not with the evidence, but with the stability of the system under variations in the input query itself. This includes semantic variations (paraphrases, synonym substitutions), typographical noise, linguistic ambiguity, and compositional perturbations that combine multiple transformations. 'Semantically, an adversarial angle,' query-level robustness is not limited to \emph{adversarial attacks} in the narrow sense---the functional requirement is a consistent, faithful response under naturally occurring user phrasing, dialectal variance, and schema mismatch. Multiple recent benchmarks collectively define this landscape, addressing retrieval and generation faithfulness under small text transformations, inherent data noise in retrieved contexts, and even counterfactual content within the knowledge base \cite{arxiv:2504.19754v1,arxiv:2503.10150v3}. Crucially, robustness must be not confused with \emph{retrieval feedback}: a system can retrieve exactly the right documents while the generator fails due to unresolved multi-hop reasoning or conflicting evidence. The combination therefore requires separate evaluation suites---benchmarks that dissociate retrieval accuracy from generation correctness to pinpoint failure sources---as highlighted by goal-oriented frameworks that distinguish \emph{generation versus retrieval ability} \cite{arxiv:2502.11371v3}.

Empirical evidence across principal leaders \cite{arxiv:2502.11371v3} suggests those integrated attribution and faithfulness frameworks---especially the \emph{generation with citation} paradigm and combined consistency boosting---produce answers that are both more verifiable and more precise, but less faithful frameworks trade off logical popularity for faithful leakage. Graph-based retrieval has demonstrated particular advantages here: by maintaining hierarchical entity- and relationship-level structure, graph-based RAG systems provide natural path-level evidence for multi-hop reasoning, strengthening both attribution-with-the-embedding route and global coherence \cite{arxiv:2502.06864v1}. However, this turns the coupling to the question of \emph{attribution in the graph}: whether the complexity cost of graph construction is justified by the attribution to other span allocated for downstream reasoning tasks.

Despite meaningful technological progress, major open challenges persist. RAG evaluation lacks a unified, information-theoretic foundation for faithfulness; existing metrics often conflate \emph{attribution quality} and \emph{generation coherence}, and attribute attribution scores on retrieved snippets without adequate per-claim grounding. The \textbf{most pressing future directions} include: (1) the development of \emph{adversarial and counterfactual attribution benchmarks} that systematically secure that systems select signal sources and reject noise---rather than simply whether they correctly cite one of K relevant documents; (2) the extension of \emph{multi-modal, hierarchical attribution}, such as when an output must cite both part of a table and a sub-image, combining text and structured evidence \cite{arxiv:2504.10074v3}; and (3) robust, adversarial attribution tests where a maliciously-compressed corpus (document with refuting information target-specific) must be corrected for---bringing the field closer to actual deployment in regulatory and safety settings. The next step will be a unified diagnosis via underlying-state probing to distinguish failure in "retrieval-selection", "retrieval-utilization", and "attribution-maintenance", to disentangle a lack of evidence from an inability to select it.


\subsection{Robustness and Security Threat Modeling}


The integration of retrieval mechanisms with generative models fundamentally expands the attack surface of modern NLP systems, introducing vulnerabilities that are distinct from both conventional language model risks and traditional information retrieval threats. While the preceding discussion established how faithfulness, attribution, and query-level robustness metrics enable us to measure truthfulness and provenance under benign and perturbed conditions, this section shifts focus from \emph{measurement} to \emph{adversarial pressure}: we now interrogate how these very properties can be deliberately undermined by malicious actors. As RAG systems are increasingly deployed in production environments where they draw upon heterogeneous, often dynamically-updated external knowledge bases, a systematic threat modeling framework is essential. The fidelity of the entire pipeline---from corpus ingestion to final generation---must be interrogated under adversarial conditions, considering the composite nature of the system as both a feature and a vulnerability \cite{arxiv:2410.12837v1,arxiv:2502.06872v1}. Drawing from the taxonomy presented in the broader surveys on robust RAG \cite{arxiv:2506.00054v1,arxiv:2502.06872v1}, we organize this analysis around four principal threat classes: data poisoning, prompt injection, structural disruption, and the emerging landscape of competitive and interference attacks---each of which directly challenges the faithfulness and attribution guarantees discussed in the preceding section.

\textbf{Data poisoning and corpus manipulation} represent the most studied class of attack in the RAG ecosystem, fundamentally subverting the attribution mechanisms that faithfulness metrics are designed to verify. The foundational vulnerability is that a single malicious document, or a small cluster of fabricated passages, embedded within a large external knowledge base can serve as a precision-guided tool to steer a generative model's output towards predetermined content. These are often called "Phantom" attacks: the attacker introduces a corpus of documents engineered to be co-retrieved alongside benign documents, thereby injecting misinformation that appears grounded in authoritative sources \cite{arxiv:2502.06872v1}. The efficacy of such attacks is contingent upon several factors, most prominently the composition of the knowledge base itself and the chunking strategies used during indexing, which directly control the semantic proximity of the malicious content to legitimate data \cite{arxiv:2506.00054v1}. Furthermore, the poisoning surface is not confined to the corpus itself; adversarial perturbations can be introduced during the retriever's fine-tuning process. By subtly manipulating the training signal (e.g., using data injection of small weighted updates with reinforcement learning), an attacker can implant a backdoor that misroutes queries for a specific topic to malicious content, compromising the entire pipeline from the retrieval stage onward \cite{arxiv:2502.06872v1}. This class of attacks directly frustrates the counterfactual attribution benchmarks identified as a future priority in the previous section, demonstrating that adversarial robustness and attribution verification are inherently intertwined.

\textbf{Prompt injection and dynamic knowledge disruption} constitute the second core attack class. RAG systems are uniquely susceptible to indirect prompt injection attacks because this retrieved knowledge, which the system treats as immutable ground truth, is placed directly into the model's context. An adversarial text can contain malicious instructions that override the model's original directives, making the generator ignore safety policies, leak data, or extract the text embedded in the prompt---which often contains the query itself \cite{arxiv:2502.06872v1}. This class also includes backdoored retrievers that return attacker-suggested results via hidden triggers, further exploiting the semantic gap between the retriever and generator components where malicious content can hide. This vulnerability directly undermines the semantic alignment and the citation integrity that attribution evaluation attempts to measure, suggesting that trustworthiness metrics must account for adversarial prompt contexts.

Given these reasonable concerns, defense mechanisms must be examined at multiple layers of the pipeline. A recent architectural framework that demonstrates advanced defense-in-depth logic is the decoupled enforcement mechanism. Rather than merely engineering the inputs, this architecture selectively discloses only the necessary policy information in the prompt---grounded in the user's query---forcing isolation from irrelevant instructions and creating a safe execution flow \cite{arxiv:2502.06872v1}. For system-level control, detection frameworks employ lightweight checks such as perplexity filtering and similarity analysis to identify remote embeddings in the retrieval stage---which act as tripwires for poisoned chunks---thereby avoiding the need for fine-tuning or full data cleaning of the entire pipeline \cite{arxiv:2502.06872v1}. By explicitly marking a source as "unretrieved" after a low-probability score and implementing a rejection process, the system remains interpretable in distinguishing its internal knowledge from the retrieved data \cite{arxiv:2502.06872v1}.

It is crucial to recognize that the threat landscape is not static but dynamically evolves as adversaries create increasingly complex attacks. A novel dimension in this space is the concept of competitive and interference attacks, where adversaries sue other malicious content in the same knowledge base to distort or neutralize a system's intended signal. In this model, an adversary's drive to manipulate the system's conclusion is counterbalanced by the possibility of another hidden adversary concurrently trying to steer the same task. Consequently, the most effective response in such a contingency differs fundamentally from the single-adversary model, and the evaluation metrics must reflect this new adversarial reality. The introduction of standardized benchmark suites, such as \emph{PoisonArena}, is a critical step forward, as these frameworks simulate a multi-source ingestior and competitive querying behavior, offering a more realistic assessment of robustness than static datasets \cite{arxiv:2502.06872v1}.

Finally, we must assess the system's resilience to \textbf{common query and composition perturbations}. These are often less targeted in nature but still affect the retriever and the generator even without explicit adversarial intent. Small lexical or semantic changes to the original query (or the query without external context) can significantly perturb the retrieval output and thus the final synthesized answer. Multi-scenario benchmarks that empirically test robustness to inherent noise, contradicting evidence, and counterfactual content are also essential \cite{arxiv:2502.06872v1,arxiv:2506.00054v1}. This directly aligns with the query-level robustness measurements described in the previous section, underscoring that robustness evaluation cannot be separated from adversarial validation.

In synthesis, the robustness of RAG demands a lifecycle, systems-level approach --- from pre-retrieval indexing decisions to end-to-end behavior --- must be continuously delivered. Adversarial threats are not a transient phenomenon but are inherently woven into the design of the entire pipeline: they challenge both the attribution metrics that measure trust and the efficiency constraints. The operational constraints and budget-aware designs that the following section will examine must incorporate adversarial resilience as a first-class requirement, rather than an afterthought. Future research must transcend a purely "defensive" perspective and embed adversarial scenarios as a baseline specification for system auditing --- in parallel with efficiency metrics. The development of causal-graph analysis and hierarchical adjudication might prove crucial to disambiguate attacks, conflicting signals, and misleading results as the complexity of RAG grows.


\subsection{Efficiency and Cost-Sensitive Performance Analysis}


Efficiency and cost-sensitive performance analysis has emerged as a critical dimension in the evaluation of Retrieval-Augmented Generation (RAG) systems, driven by the practical realities of deployment where operational constraints often dictate architectural choices. The economic feasibility of RAG hinges on carefully balancing the computational expenditures of both retrieval and generation against the quality of the final output, particularly because the augmentation process introduces substantial overhead in terms of latency, throughput, and memory footprint beyond the already significant costs of underlying large language models (LLMs). The retrieval stage, typically involving nearest neighbor search over densely encoded billion-scale vector indexes, presents a major bottleneck, while the generation stage incurs quadratic attention costs inflated by the injection of even modest numbers of retrieved evidence documents. From a system perspective, the end-to-end budget spans not only compute and storage for indexing, but also the implied costs of query transformation, scoring / re-ranking, and the decoding-time interventions used to ensure fidelity \cite{arxiv:2506.00054v1}.

The core of efficiency-aware RAG analysis focuses on latency, throughput, and memory characterization. Empirical findings highlight that performance on both short-form and multi-hop question answering tasks improves as the number of retrieved contexts increases, yet such gains plateau---and may even decline---beyond an ideal threshold, suggesting a nuanced balance between retrieved quantity and generation load \cite{arxiv:2502.14759v1}. In practice, the injection of large context windows can lead to "attention dilution," degrading model performance and escalating computational cost. To address this, the research community has proposed strategies for context compression, including token-level pruning, hierarchical summarization, and latent semantic decoding, all of which aim to reduce the volume of data sent to the generator while preserving critical information. In contrast, approaches like cache-augmented generation (CAG) completely circumvent the retrieval stage by pre-computing and caching relevant knowledge, trading a higher storage footprint for significantly lower online query latency \cite{arxiv:2506.06704v1}. Complementing these algorithmic strategies are hardware- and system-level accelerations, such as memory-centric architectures and quantization for approximate nearest neighbor search, which collectively enable high-throughput operation with manageable operational footprints.

A key dimension for cost-sensitive analysis is the strategic selection of an optimal operating point via automatic component selection. Frameworks such as RAGe and VERA have been developed to automatically configure RAG pipelines to either meet specific accuracy targets or minimize end-to-end runtime by assessing quality-latency trade-offs across candidate retrievers and generators \cite{arxiv:2506.00054v1}. Complementing these configuration-level decisions, lightweight search agents trained via reinforcement learning--using rewards like the "Gain Beyond RAG" that measures utility improvement over non-retrieval baselines--can dynamically decide whether to perform retrieval at all, thereby significantly cutting communication and latency costs in federated or multi-source settings \cite{arxiv:2505.14146v2,arxiv:2502.19280v2}. Multi-armed bandit approaches extend this to adaptive retrieval strategy selection, balancing exploration and exploitation to reduce overall retrieval steps while maintaining answer accuracy \cite{arxiv:2412.01572v4}. At an architectural level, parameter-based RAG methods such as Parametric RAG and Dynamic Parametric RAG present a fundamental re-alignment by embedding knowledge directly into model parameters to reduce online computational costs and the need for context integration \cite{arxiv:2501.15915v1,arxiv:2503.23895v4}. Similarly, the token-efficient TeaRAG framework demonstrates that by extracting and focusing on a knowledge association graph, output tokens can be reduced by up to 61\%, yielding a substantial reduction in memory and latency with minimal impact on answer accuracy \cite{arxiv:2511.05385v2}.

Taken together, these metrics and optimization frameworks reveal a complex and often contradictory landscape where improving the efficiency of one component can harm the system overall. Advancing the field therefore requires a "system-guided" approach: parameterizing and jointly optimizing the end-to-end pipeline, including corpus scaling, index structures, hardware accelerators, and model compression \cite{arxiv:2506.00054v1}. Furthermore, a system-level viewpoint that abstracts exact model choices and treats both retrieval and generation as flexible modules will be crucial for constructing latency- and cost-constrained RAG systems that are not only practical but also theoretically sound. Future research must reconcile the goal of maximum accuracy with the hard limitations of real-world computational resources using the development of principled, budget-aware designs.


\section{Applications, Frontiers, and Future Horizons}



\subsection{Domain-Specific Deployments Across Knowledge-Intensive Fields}


The translation of retrieval-augmented generation from research prototypes to production-grade systems is perhaps most tangibly realized in knowledge-intensive fields where the cost of factual error is prohibitive and the demand for verifiable, current, and domain-specific grounding is non-negotiable. The generic RAG architecture---comprising a retriever over an external corpus and a generator conditioned on the retrieved evidence---is not universally deployable without careful, domain-conscious adaptation. The pragmatic deployment in specialized sectors necessitates a systematic re-engineering of the RAG pipeline across three critical axes: (1) the structuring of the knowledge base to align with domain-specific data models and vocabularies, (2) the design of retrieval strategies to handle heterogeneous and often semi-structured data, and (3) the fine-tuning or prompting of the generator to adhere to domain-specific operational constraints and normative reasoning \cite{arxiv:2501.13958v3,arxiv:2312.10997v5}. These field-specific configurations reveal fundamental design trade-offs between the flexibility of generic semantic search and the precision of structured or rule-based domain logic.

In the healthcare and biomedical sector, the stakes of factual recall are exceptionally high, and the architecture of RAG systems reflects this. Deployments must not only address the semantic gap between layperson or even specialist queries and dense clinical vocabulary but also navigate strict privacy constraints that preclude sending protected health information to public index infrastructure. Beyond mere vocabulary alignment, there is a pressing need for decision support that synthesizes across a heterogeneous, multi-modal corpus that includes clinical guidelines, electronic health records (EHRs), and the ever-evolving peer-reviewed literature \cite{arxiv:2312.10997v5,arxiv:2405.06211v3}. The critical technical challenge within this field is that of \textbf{structured authority exploitation}: it is crucial to ensure that primary clinical contexts and patient context take precedence over the model's parametric priors. While vector-based chunk retrieval (e.g., a PubMed-derived index) retains importance, researchers have increasingly turned to \textbf{graph-based RAG} to explicitly model relationships and bridge the context gap, e.g., embedding patients' specific demographics and drug interactions as nodes \cite{arxiv:2502.06864v1,arxiv:2502.01113v3}. This move towards a logical, graph-centric index and \textbf{corrective retrieval} mechanisms---where a pre-retrieval evaluator is used to filter or reject low-relevance, potentially misleading evidence before it reaches the generator---proves essential for mitigating risks \cite{arxiv:2401.15884v2,arxiv:2410.07176v2}. The future of clinical RAG lies in effectively modeling temporal dependencies and reasoning across longitudinal, patient-specific episodic memory and current bio-medical guidelines in a single, sequential, and evidence-credentialed step.

Similarly, high-stakes regulation and transactional fields such as law and finance place unprecedented emphasis on the \textbf{temporal integration} and the \textbf{semantic framing} of evidence. These domains are volatile and heavily structured, with legal and financial decisions resting heavily on document versions, explicit dates, regulatory updates, and hierarchical contract clauses that cannot be compromised. The application of a na\"{i}ve chunk-level retriever over a monolithic corpus is fundamentally flawed. For example, a legal statute may have been amended; a financial technical report may be superseded by a more recent filing; the context between them was missing. In these cases, the nature of the fact goes beyond the difference in the text itself; the ability to unify into a coherent, interconnected, and chronology-aware picture is essential. The architecture shifts toward integrating \textbf{knowledge graph-extended retrieval}, injecting a \textbf{structural} knowledge graph into the pipeline \cite{arxiv:2502.06864v1,arxiv:2502.06864v1}. This allows the system to handle the retrieval of information in a structured, semantic graph schema; a dedicated retrieval mechanism can correctly enforce temporal context ordering and traverse structural relationships---by exploring entity relationships in these graphs (e.g., with graph algorithms) to model logical semantics, contracts, and market volatility \cite{arxiv:2501.13958v3}. In finance, this is also a form of stress relief; the \textbf{design objective} is to shift the balance and to perform the \textbf{continuous and incremental indexing} of the stock data, ensuring the retrieval engine doesn't become stale in a matter of minutes, not just timely and quality data \cite{arxiv:2405.06211v3}. This has led to the emergence of \textbf{cache-augmented generation (CAG)} for high temporal information; when the knowledge base is manageable and bounded (a specific company's clauses, a quarterly report filings), preloading the entire corpus into the LLM's extended context and deriving its runtime parameters eliminates retrieval latency and, more importantly, reduces the risk of retrieving an outdated or conflicting version. For such operational use, it should not be necessary to sacrifice computational resources \cite{arxiv:2412.15605v2}.

Beyond regulated fields, \textbf{enterprise and customer service deployments} present a different yet equally stringent set of constraints, centered around user-centric knowledge grounding, combined language, and high-volume concurrency. RAG systems here are not just answering questions; they are engendering multi-turn, memory-based recipe. In customer-oriented agent interactions, an knowledge acquisition can be addressed through grammar and domain-specific vocabulary, but the principles behind this best practice have moved towards a \textbf{compositional, multi-source retrieval} method rather than a monolithic index. The challenge for entity resolution becomes one of \emph{source decomposition and synthesis}: instead of injecting the entire set of retrieved chunks, the RAG engine must orchestrate a decomposition-to-recomposition framework and extract a core set of evidence. This is like in the case of an online technical support system, where a "correct" answer to the \textbf{customer-centric question} is useless if not grounded in a strict set of documentation; and the combination of document context in order to answer an resource query is required. The system's capability to reduce system--level prompt bloat via an explicit retrieval-graded and augmented decision boundary is akin to what information retrieval and generation in agents is proposed in \cite{arxiv:2505.03275v1}. Furthermore, a new set of \textbf{enterprise's related fine-tuning}: adding to the standard ICL architecture with \textbf{retrieval-robust fine-tuning} can be done; by giving the generation model a task to get a larger, but not limited, domain coverage and to cope with the many incoming issues, real-world noise \cite{arxiv:2501.18365v1}. Fine-tuning based on generated (curated) databases that artificial code and the duel. This often includes \textbf{retrieval-aware data-augmentation}, generating a synthetic QA corpora for the tasks which captures knowledge of the voice and local terms in the field, directly and in training the generator. The approach improves the alignment semantics between the search layer and the end-user dynamics of language.

Along with these domain-centric architectures, the last key point is the risk-aware robustness of the knowledge process: RAG systems are only as critical as the reliability of their components. The real deployment shows a shift from reliability to \textbf{reliability at operation}: a commitment to instrumenting failure modes that are specific to the field \cite{arxiv:2401.05856v1}. In the medical domain, a legitimate treatment at least as confident as the reliability of each intermediate retrieval step; in finance, a \textbf{counterfactual robustness} claim risks compliance penalties; in enterprise, a repeated failure and harsh recovery erodes user trust. A clear observation across these domains indicates that the best effective RAG systems are not monolithic, but extensible and carefully \textbf{scaffolded}. They rely on a graph-augmented approach to handling document situations across edge cases. This points towards the top-tier deployment of RAG as a visible, adaptive and demanding model of evolution from a bounded pipeline; this supports \textbf{agentic orchestration} and the strategic node: the LLM can pro-actively reason about decidable criteria for retrieval, deciding when to retrieve to support the decision, and selecting among diverse tools and data sources specific to the target domain \cite{arxiv:2501.09136v4,arxiv:2411.19443v1}. For domains where reasoning and decision-making are critical (e.g., first, do no harm; \emph{stare decisis}; Net Realization), the frontier is \emph{not} in the parametric model's cognitive limitations but in the use of a logic and ecosystem such as a prober to make the system sharper in error detection or evidence triage. The focus must remain on closing the capability horizon for specific user tasks through a \emph{fair and auditable} process rather than from "an ungrounded". Future RAG work in knowledge systems must embrace a holistic memory design, where generation to serve users within a domain converges on the RAG architecture: transparent---to enable fine-grained control and safe guardrails. For, after all, the clinical and domain-specific RAG is not just about answering, but about answering with justifiable \emph{source-grounded and factually accurate} reasoning.


\subsection{Multi-Modal and Structured Knowledge Integration}


The evolution of retrieval-augmented generation beyond purely textual corpora represents a critical frontier for grounding large language models in the full richness of human knowledge, directly extending the domain-specific and structurally complex systems discussed above. This expansion is driven by the recognition that real-world documents and domain-specific repositories are inherently multi-modal and structurally complex, containing images, tables, audio, and explicit relational links that are lost in flat text representations \cite{arxiv:2405.06211v3}. The integration of these diverse sources profoundly alters the design of the retrieval component and the augmentation strategy, demanding a move from simple token-based matching toward semantic alignment across heterogeneous representational spaces---a shift that sets the stage for the agentic orchestration architectures explored in the next section.

For multi-modal content, the central design question is at what point in the pipeline visual, audio, and textual knowledge should be fused. The most straightforward approach is to reduce all modalities to a textual representation via a multimodal large language model, generating captions or textual summaries of images or audio files to populate a standard vector index. This yields a corpus that is compatible with established retrieval and augmentation pipelines, leveraging the semantic comprehension of the vision-language model to capture context that would be lost in a simple OCR or ASR transcription. However, this early fusion is prone to information loss, as the richness and specific details of the source---such as the precise spatial layout of a document or the emotional tone of a recording---are discarded during the summarization step \cite{arxiv:2410.10594v2}. To preserve this fidelity, a more advanced approach embeds the documents directly as image embeddings with a multimodal encoder, treating raw pages as the atomic retrieval unit \cite{arxiv:2502.11371v3}. This method retains all visual information, layout, and images, avoiding the semantic gap of parsing, and empirically demonstrates superior performance in both retrieval and generation---achieving significant end-to-end improvement over text-based RAG pipelines \cite{arxiv:2502.11371v3}. The emerging frontier is a hybrid strategy that selects the representation based on query type, maintaining both raw modalities and their semantic summaries in a unified index, but this remains largely unexplored.

Paralleling this exploration along the structural dimension of the system is research into knowledge graph integration and hierarchical structuring, which further compound the heterogeneity of the knowledge landscape. Graph-based RAG systems address the inherent limitations of flat, chunk-based retrieval---which struggles with complex inter-dependencies and contextual awareness---by explicitly modeling relationships and hierarchies across documents \cite{arxiv:2510.02657v3}. This enables multi-hop reasoning and more robust retrieval of related entities, unlike flat chunking which tends to isolate fragments of context. However, both paradigms suffer from substantial overhead, index noise, and structural incompleteness. Graph-based retrievers are limited by the structural clarity of the underlying knowledge, and they struggle to manage dynamic corpora and multi-source integration \cite{arxiv:2503.04338v2,arxiv:2502.11371v3}. Targeted techniques such as flow-based pruning of redundant relational edges for efficient graph traversal, single-step generation or incremental-update-oriented designs have been explored to reduce this noise and support evolving data landscapes \cite{arxiv:2502.11371v3,arxiv:2208.07652v1,arxiv:2502.11371v3}.

Given these concrete matured regimes, the deployment bottleneck becomes how the system orchestrates this diversity---linking text with graphs, and parametric knowledge with non-parametric memory \cite{arxiv:2605.15156v2,arxiv:2401.14887v3}. True innovation is a pipeline that makes retrieval cross-modal and multi-source (hybrid) deliberately, e.g., by fusing dense embeddings from visual content with regression to answer a single query \cite{arxiv:2502.00592v2}. This involves overcoming several challenges: the differing modalities feature semantics and corpus distribution, the chasm between discrete entity links vs. the continuous nature of visual embeddings, and the inherent tradeoff of graph-size explosion vs. retrieval precision.

Looking forward---and directly anticipating the discussion of multi-agent orchestration---a promising avenue is the integration of LLM agents with persistent structured knowledge memory, where hierarchical chains and graph-based abstraction layers become formalized, progressively enriched, and able to support continuous adaptation \cite{arxiv:2412.05838v1}. This would allow future RAG systems to achieve evolving intelligence through unifying---rather than juxtaposing---the world's structured, dynamic, and embodied knowledge, illustrate that grounding generative models on heterogeneous or broadly heterogeneous and contextually integrated knowledge remains the core future objective \cite{arxiv:2605.15156v2}.


\subsection{Agentic, Multi-Agent, and Tool-Augmented Architectures}


The evolution from single-path retrieval-augmented generation (RAG) pipelines toward agentic, multi-agent, and tool-augmented architectures represents a fundamental shift in how external knowledge is integrated with large language model (LLM) reasoning. Rather than treating retrieval as a discrete, one-time operation, these systems reconceptualize the LLM as an autonomous orchestrator capable of multi-step planning, dynamic tool selection, iterative feedback, and self-correction. This transition enables RAG systems to address increasingly complex, multi-hop, and reasoning-intensive tasks that demand adaptive and interactive knowledge acquisition.

At the core of this shift is the emergence of the LLM as an agentic controller that transcends the roles of simple query formulator and answer synthesizer. In these designs, the model makes autonomous decisions about \emph{when} retrieval is necessary, \emph{what} information to seek, and \emph{how} to integrate the obtained knowledge into an evolving reasoning process. For instance, the TeaRAG framework introduces a token-efficient agentic RAG model in which the LLM not only performs multi-round reason-retrieval iterations---each round guided by its own hidden states---but is also explicitly incentivized to regulate its own reasoning process and to balance informational sufficiency against computational cost \cite{arxiv:2511.05385v2}. Instead of a single retrieval call followed by a fixed generation step, the agentic framework embeds the retrieval decision \emph{within} the generation loop itself, enabling a more organic interleaving of evidence collection and answer formulation. This marks a clear departure from earlier frameworks in which the LLM statically consumed a pre-retrieved context.

A significant trend within this paradigm is the \textbf{hierarchical, intent-aware decomposition of retrieval tools}. Agentic systems increasingly implement layered retrieval interfaces that range from a single-pass semantic search to more complex, multi-stage processes. The practice of dynamic tool selection across a spectrum of retrieval mechanisms represents a critical design choice that directly influences the system's reasoning capacity and the quality of grounding. The GeAR system extends this concept by integrating a graph-based retrieval mechanism into an agent framework with a "graph expansion" module and an adaptive agent loop, allowing for both independent and iterative retrieval steps over a dynamic textual knowledge base---a capability that proves especially beneficial for multi-hop QA tasks \cite{arxiv:2412.18431v2}. Similarly, the AGRAG framework addresses a key limitation of graph-based RAG by substituting unreliable LLM entity extraction with a statistics-based method and frames the graph reasoning process as a subgraph generation problem, the Minimum Cost Maximum Influence (MCMI) problem, thereby enabling explicit, noise-reduced reasoning paths that guide the model's attention toward answer-relevant content \cite{arxiv:2511.05549v2}. These contributions indicate a shift from data retrieval to knowledge engineering in which the agent itself---through its decision-making abilities over an extended graphical structure---actively mitigates the inherent noise and incompleteness of the underlying knowledge corpus.

The \textbf{multi-agent orchestration} paradigm further creatively advances the capabilities of agentic RAG by distributing roles among specialized agents. This distributed cognitive architecture demonstrates that a single LLM, even with sophisticated retrieval strategies, may face bottlenecks when processing diverse data sources due to the pressure on a single context window and the difficulty of switching between tasks. A collaborative multi-agent RAG system pinpoints a single point of failure: the ability of a single agent to handle trade-offs across multiple retrieval contexts and to generate from diverse data---not just text document stores, but relational, NoSQL, and graph databases---without performance degradation. This multi-agent system divides the query generation work among specialized agents, each optimized for a specific data source, and in doing so reduces token consumption, increases resilience, and maintains accuracy even at scale \cite{arxiv:2412.05838v1}. Additionally, despite performance gains from parallelism in this multi-agent model, moving from a single to a multi-agent design introduces new encapsulation complexity---how to program inter-agent communication, coordinate retrieval efforts, and aggregate results into a single coherent final answer remains a formidable open challenge.

Crucially, agentic systems do not only decide \emph{when} to retrieve but also \textbf{how many and in what sequence to process a tool such as a retriever or a reasoning step}. Current advances such as TeaRAG demonstrate an \textbf{efficient-mode} approach to optimizing the learning process---the framework uses a knowledge-based reward that encourages the minimization of long deliberations, punishes unnecessary reasoning steps, and compresses the output token count at inference time \cite{arxiv:2511.05385v2}. These designs turn a system optimization constraint into a model behavior constraint. It does so by inserting a learning objective into the agent's decision-making, which aligns its efficiency--accuracy trade-offs with the \emph{per-token cost} of the inference process, rather than carefully optimizing it as a post-hoc system-level task.

Looking to the future, the extension of agentic approaches from text-only retrieval to multi-modal sources remains a core evolutionary need. Although the \textbf{multi-agent architecture} remains strong, most existing frameworks are directed towards textual sources. With the ever-growing prevalence of multimodal documents, extending these agentic protocols to \textbf{vision- or audio-based RAG} systems enhancing AI grounding with multimodal data is both a logical and necessary progression \cite{arxiv:2410.10594v2,arxiv:2502.14727v1}. Critically, this architectural evolution shares with the concept that retrieval is an \textbf{action}, not a black-box pre-processing stage. By embedding or integrating these decisions with generation and offline training, future RAG systems may demonstrate significantly more robust, self-diagnosing, and efficient grounding behaviors, reasoning increasingly closer to the model's output.


\subsection{Long-Horizon Memory, Personalization, and Continuous Learning}


The evolution from single-query, stateless pipelines toward systems endowed with persistent, episodic, and lifelong memory structures represents a paradigm shift that builds directly upon the agentic and tool-augmented architectures discussed in the previous section. Whereas those systems reconceptualized the LLM as an orchestrator capable of dynamic tool selection and iterative feedback within a single task, the emerging class of memory-augmented systems extends this principle across temporal horizons. These frameworks seek to unify short-term and long-term episodic memory with user profile modeling, thereby enabling coherent behavior across extended interactions \cite{arxiv:2409.19401v1}. This transition is not merely an incremental improvement but a fundamental reconceptualization of the system's relationship with its own history and with the individual user it serves---an evolution that lays the groundwork for the long-horizon robustness and self-adaptation requirements that will be examined in the following section.

Central to this evolution is the design of explicit memory modules that persist beyond the immediate inference context. Editable Memory Graphs (EMG) exemplify this approach: user data---continuously generated on personal devices---is organized into an editable graph structure that a personalized agent queries via RAG techniques. Crucially, the editability and selectability of this memory graph are optimized using reinforcement learning, addressing the threefold challenge of data collection, dynamic updating, and relevance selection. Empirical validation on real-world datasets demonstrates an approximate 10\% improvement over the best existing approaches and successful deployment within actual smartphone AI assistants, confirming that user-specific grounding can substantially enhance downstream task performance \cite{arxiv:2409.19401v1}. Similarly, for finer-grained conversational memory, the GAM-RAG framework introduces a training-free, gain-adaptive memory that accumulates retrieval experience from recurring or related queries. GAM-RAG constructs a lightweight, relation-free hierarchical index whose links capture latent co-occurrence, and it updates sentence-level memories during inference via an uncertainty-aware, Kalman-inspired gain rule that balances fast adaptation for novel signals with conservative refinement for stable ones. This yields a 3.95\% performance improvement over the strongest baseline, with a striking 61\% reduction in inference cost, demonstrating that dynamic memory can confer not only higher accuracy but also substantial efficiency gains \cite{arxiv:2603.01783v1}.

Despite these advances, a fundamental dimension of long-horizon memory remains comparatively underdeveloped: personalization. Current RAG systems generally treat documents and external corpora as globally relevant units of knowledge \cite{arxiv:2503.10677v3}. However, the discriminating utility of a retrieved fact is profoundly conditional on the user's long-term preferences and immediate linguistic context. Statistical regularities in the user's behavioral trajectory---while powerful---are often insufficient, as they ignore the latent drift in preferences over time. A more promising direction is the design of user-centric grounding, where retrieval is conditioned not solely on the query but also on a user profile that itself co-evolves with interaction history. In this vein, the multi-agent ARAG framework moves beyond naive personalization: a User Understanding Agent summarizes long-term and session-specific behavior, an NLI Agent evaluates semantic alignment between retrieved evidence and inferred intent, and an Item Ranker Agent produces contextualized ranked outputs. The reported improvements---up to 42.1\% in NDCG@5 \cite{arxiv:2506.21931v2}---substantiate the hypothesis that treating personalization as a reasoning- and retrieval-integrated problem, rather than a static filter, yields qualitatively better grounding. Yet a key open research question is how to train augmenting modules for such user-centric RAG systems without degrading the rich capabilities of the base LLM, particularly in continual-learning settings where preferences shift gradually.

A third critical axis concerns the scalability of memory management under sustained deployment. Long-horizon reasoning requires models to attend to both recent and distant episodes, yet naive context concatenation quickly exceeds computational budgets and incurs attention dilution. Memory-centric RAG systems must therefore be designed with multi-agent techniques to decompose cognitive load across specialized contributors, as exemplified by the division of labor in ARAG \cite{arxiv:2506.21931v2} and the teacher-student information refinement in MAIN-RAG \cite{arxiv:2501.00332v1}. Additionally, when such systems are fine-tuned at deployment, they risk catastrophic forgetting of the model's general competencies---a possibility, though discussed in the context of long-horizon personalization, ultimately links back to the issue of balancing specialized adaptation with robustness that surfaced in the multi-agent orchestration of the previous section. SelfAug directly addresses this failure-mode in a RAG context by aligning the fine-tuned model's distribution with its original one via self-distribution alignment on input sequence logits, and demonstrates that this not only mitigates forgetting but also improves downstream RAG accuracy \cite{arxiv:2509.03934v1}. This suggests that continual learning of a personalized model can co-exist with general capability preservation, but it requires training and memory mechanisms that remain cognizant of distributional drift.

From a broader perspective, the pursuit of long-horizon memory should be decoupled from the problem of retrieval-augmented generation itself. The objective is not simply to retrieve more, but to decode and maintain a model of the user over time---a personalization that extends beyond data semantics to a memory model informed by all previous interactions, including user-specified memories and multi-step usage patterns. Current systems are beginning to consolidate user memories in editable graph topology and to automate memory collection through reinforcement-learned summarization \cite{arxiv:2409.19401v1}. Future frameworks must further incorporate theories from schema-based learning in cognitive neuroscience, as pioneered by GAM-RAG, to design memory modules that anticipate, rather than merely react to, recurring information needs \cite{arxiv:2603.01783v1}. Finally, as systems reach toward lifelong operation, it will be necessary to revisit the fundamental trade-off between cognitive persistence and degradation from entropy accumulation in the internal representations that encode episodic memory. The integration of latent representation and continual, enabling incremental indexing will be essential to make these long-horizon, deeply personalized agents both persistent and stable at a global deployment scale---a goal that aligns with the broader move toward test-time self-adaptation and robust multi-turn evaluation discussed in the following section. Achieving this unification of memory, personalization, and continual adaptability will see RAG evolve from a strong but narrow tool into a truly long-term, user-centric AI layer.


\subsection{Emerging Frontiers, Evaluation, and Open Challenges}


As retrieval-augmented generation matures from a pipeline engineering discipline into a more scientifically grounded field, several fundamental questions emerge that will define its next evolution. The most pressing theoretical inquiry concerns the very conditions under which augmentation provides measurable utility. While RAG systems demonstrably enhance factual accuracy and mitigate hallucination, the field lacks principled models explaining when retrieval---with its attendant latency, computational overhead, and potential for introducing noise---justifies its cost relative to alternative approaches such as long-context models \cite{arxiv:2312.10997v5}. Hybrid routing mechanisms that dynamically select between direct inference and retrieval-based processing have shown promise, particularly through syntax-aware complexity analysis that can classify queries and route accordingly \cite{arxiv:2602.03578v1}. Similarly, multi-armed bandit approaches treating retrieval strategies as arms selected dynamically demonstrate that such adaptation can yield state-of-the-art performance while reducing retrieval costs \cite{arxiv:2412.01572v4}. However, these approaches remain largely heuristic; formal models characterizing the marginal value of augmentation as a function of task difficulty, internal knowledge confidence, and corpus characteristics remain a critical open frontier.

The evaluation of RAG systems is likewise entering a period of introspection. Traditional component-wise metrics---separately measuring retrieval precision and generation quality---often fail to capture systemic interactions, leading to a misalignment between benchmark scores and deployed performance \cite{arxiv:2506.00054v1}. This has motivated unified evaluation frameworks such as RAGAs, which provides reference-free metrics spanning answer faithfulness and context utilization \cite{arxiv:2309.15217v1}. Similarly, taxonomy-driven error analysis has proven essential, revealing that a substantial portion of RAG failures stem from systemic issues in retrieval defects and their cascading effects into generation \cite{arxiv:2510.13975v2}. Multi-hop scenarios complicate evaluation further: benchmark design that explicitly tests retrieval and reasoning over multiple, distributed evidence pieces reveals that both state-of-the-art retrievers and LLMs perform unsatisfactorily on such queries \cite{arxiv:2401.15391v1}. Understanding internal knowledge processing---such as how models weight parametric knowledge against retrieved context via causal mediation analysis---has revealed that models, when both are available, rely on the context over parametric knowledge, though this pattern requires careful controls and its generality remains open \cite{arxiv:2410.05162v1}. This finding underscores the need for a new generation of evaluation protocols designed to systematically disentangle conflicting dominant pipeline signals from final-output variation factors.

A deeper theoretical understanding of knowledge conflicts is another research frontier with field-wide implications. Recent results show that post-retrieval generation suffers from over-reliance on internal parametric knowledge stored in particular FFN layers \cite{arxiv:2502.15543v4}, and that forced fidelity suppression of parametric knowledge, while increasing context adherence, can actually undermine a model's ability to correctly interpret evidence \cite{arxiv:2506.08938v2}. This synergistic yet adversarial duality---between parametric memory and nonparametric evidence---requires formal treatment through conflict resolution mechanisms. Graph-based conflict resolution has shown potential by enabling fact-level triple extraction and entropy-based filtering \cite{arxiv:2511.10375v1}, while other frameworks for transparent conflict resolution combine dual contrastive encoders with SNR-based weighting to make the model's knowledge-determination process observable and controllable \cite{arxiv:2601.06842v1}. Resilient approaches have shown promise in the worst-case scenario by eliciting internal knowledge and consolidating external with source-awareness \cite{arxiv:2410.07176v2}. Finally, as RAG systems jointly tune their components, a holistic assessment of learning dynamics in multi-agent and multi-objective settings will be essential for advancing the field.

More sophisticated evaluation protocols are crucial for the rank domain/vertical/vertical domains. These frameworks need to move beyond the static one-shot question-and-answer and single-turn tasks to evaluate multi-turn robustness, the impact of degraded or poisoned corpora, and the system's overall utility over extended temporal horizons \cite{arxiv:2401.05856v1}. Given that software engineers now know that validation of RAG is only feasible through systematic, end-user operation, test-time self-adaptation without re-training will be decisive \cite{arxiv:2401.05856v1}. The design of adaptive and self-improving systems---which dynamically gather and learn from feedback, refine the sufficiency of their internal representations, and employ self-practicing or distillation-based mechanisms to improve inference-state awareness---offers a research vector toward practical reliability \cite{arxiv:2505.02811v2}. By developing a robust and reflective system that establishes conditions under which retrieval becomes viable through principled integration with long-context and verification methods, we can overcome the intrinsic knowledge bottlenecks that are difficult to address otherwise.


\section{Conclusion}


This survey has traversed the rapidly evolving landscape of Retrieval-Augmented Generation, tracing its trajectory from foundational paradigms---exemplified by early work on differentiable access to non-parametric memory \cite{arxiv:2005.11401v4}---to the sophisticated, adaptive, and agentic systems that define the current frontier \cite{arxiv:2501.09136v4,arxiv:2506.10408v1}. Our synthesis reveals a field characterized by a profound architectural diversification, moving from monolithic pipeline designs toward a rich design space that includes pre-training-based integration, modular cascades, and fully unified models \cite{arxiv:2410.03439v3}. A central theme is the growing sophistication of system workflows, which have evolved from single-step retrieve-then-generate processes to dynamic, multi-agent frameworks that can plan, reflect, and iteratively refine their knowledge acquisition strategies \cite{arxiv:2502.18635v2,arxiv:2305.06983v2,arxiv:2411.19443v1}. These systems are increasingly capable of dynamically negotiating the balance between their internal parametric knowledge and external, non-parametric evidence, a capability highlighted by the development of adaptive retrieval mechanisms that trigger external access only when necessary \cite{arxiv:2212.10511v4,arxiv:2410.01782v1,arxiv:2502.01142v2}. The integration of graph-based knowledge structures represents another significant advancement, providing the relational and hierarchical context vital for complex, multi-hop reasoning \cite{arxiv:2410.05779v3,arxiv:2502.06864v1,arxiv:2501.13958v3}. Furthermore, the field has expanded beyond text to embrace multi-modality, with systems designed to directly process and retrieve information from images and other non-textual sources \cite{arxiv:2410.10594v2,arxiv:2503.18016v1}.

Methodologically, the survey has highlighted the transformative impact of retrieval on the generation process itself. The integration of external knowledge has moved beyond simple prompt concatenation to encompass parameter-efficient fine-tuning, decoding-time interventions, and the development of synergistic training objectives where the retriever and the generator are jointly optimized to achieve a shared goal \cite{arxiv:2110.07752v2}. This necessitates a critical evaluation of system components, where limitations such as the inherent imperfections of retrieval \cite{arxiv:2401.14887v3} and the complex resolution of conflicts between internal parametric memory and external evidence \cite{arxiv:2410.07176v2} necessitate sophisticated solutions. Despite the impressive capabilities and comprehensive evolutions, substantial open challenges remain. These include ensuring robust performance in the face of adversarial manipulations like corpus poisoning or prompt injection \cite{arxiv:2509.23519v2,arxiv:2401.05856v1}, and maintaining efficacy in continually evolving knowledge environments without constant, resource-intensive re-training \cite{arxiv:2505.17005v1,arxiv:2502.14802v2}. While long-context windows present a promising alternative to external retrieval, they fail to eliminate the need for RAG, as their performance degrades with a vast number of inputs, including "hard negatives" \cite{arxiv:2410.05983v1,arxiv:2401.05856v1}. Therefore, our comprehensive analysis positions RAG not merely as a temporary patch for the inherent parametric limitations, but as a core architectural strategy for building more reliable, reliable, and adaptable language systems \cite{arxiv:2403.03187v1} evidence from the bibliography and applications \cite{arxiv:2410.15944v1}.

\begin{thebibliography}{128}
\setlength{\itemsep}{2pt}

\bibitem{arxiv:2312.10997v5}
Yunfan Gao, Yun Xiong, Xinyu Gao, Kangxiang Jia, Jinliu Pan, Yuxi Bi \emph{et al.}.
\newblock \emph{Retrieval-Augmented Generation for Large Language Models A Survey}.
\newblock arXiv:2312.10997v5, 2023.
\newblock \url{https://arxiv.org/abs/2312.10997v5}

\bibitem{arxiv:2410.12837v1}
Shailja Gupta, Rajesh Ranjan, Surya Narayan Singh.
\newblock \emph{A Comprehensive Survey of Retrieval-Augmented Generation (RAG): Evolution, Current Landscape and Future Directions}.
\newblock arXiv:2410.12837v1, 2024.
\newblock \url{https://arxiv.org/abs/2410.12837v1}

\bibitem{arxiv:2212.10511v4}
Alex Mallen, Akari Asai, Victor Zhong, Rajarshi Das, Daniel Khashabi, Hannaneh Hajishirzi.
\newblock \emph{When Not to Trust Language Models Investigating Effectiveness of Parametric and Non-Parametric Memories}.
\newblock arXiv:2212.10511v4, 2022.
\newblock \url{https://arxiv.org/abs/2212.10511v4}

\bibitem{arxiv:2506.00054v1}
Chaitanya Sharma.
\newblock \emph{Retrieval-Augmented Generation: A Comprehensive Survey of Architectures, Enhancements, and Robustness Frontiers}.
\newblock arXiv:2506.00054v1, 2025.
\newblock \url{https://arxiv.org/abs/2506.00054v1}

\bibitem{arxiv:2005.11401v4}
Patrick Lewis, Ethan Perez, Aleksandra Piktus, Fabio Petroni, Vladimir Karpukhin, Naman Goyal \emph{et al.}.
\newblock \emph{Retrieval-Augmented Generation for Knowledge-Intensive NLP Tasks}.
\newblock arXiv:2005.11401v4, 2020.
\newblock \url{https://arxiv.org/abs/2005.11401v4}

\bibitem{arxiv:2305.06983v2}
Zhengbao Jiang, Frank F. Xu, Luyu Gao, Zhiqing Sun, Qian Liu, Jane Dwivedi-Yu \emph{et al.}.
\newblock \emph{Active Retrieval Augmented Generation}.
\newblock arXiv:2305.06983v2, 2023.
\newblock \url{https://arxiv.org/abs/2305.06983v2}

\bibitem{arxiv:2403.03187v1}
Akari Asai, Zexuan Zhong, Danqi Chen, Pang Wei Koh, Luke Zettlemoyer, Hannaneh Hajishirzi \emph{et al.}.
\newblock \emph{Reliable, Adaptable, and Attributable Language Models with Retrieval}.
\newblock arXiv:2403.03187v1, 2024.
\newblock \url{https://arxiv.org/abs/2403.03187v1}

\bibitem{arxiv:2405.07437v2}
Hao Yu, Aoran Gan, Kai Zhang, Shiwei Tong, Qi Liu, Zhaofeng Liu.
\newblock \emph{Evaluation of Retrieval-Augmented Generation: A Survey}.
\newblock arXiv:2405.07437v2, 2024.
\newblock \url{https://arxiv.org/abs/2405.07437v2}

\bibitem{arxiv:1809.05296v5}
Deng Cai, Yan Wang, Victoria Bi, Zhaopeng Tu, Xiaojiang Liu, Wai Lam \emph{et al.}.
\newblock \emph{Skeleton-to-Response Dialogue Generation Guided by Retrieval Memory}.
\newblock arXiv:1809.05296v5, 2018.
\newblock \url{https://arxiv.org/abs/1809.05296v5}

\bibitem{arxiv:2202.01110v2}
Huayang Li, Yixuan Su, Deng Cai, Yan Wang, Lemao Liu.
\newblock \emph{A Survey on Retrieval-Augmented Text Generation}.
\newblock arXiv:2202.01110v2, 2022.
\newblock \url{https://arxiv.org/abs/2202.01110v2}

\bibitem{arxiv:2501.09136v4}
Aditi Singh, Abul Ehtesham, Saket Kumar, Tala Talaei Khoei, Athanasios V. Vasilakos.
\newblock \emph{Agentic Retrieval-Augmented Generation: A Survey on Agentic RAG}.
\newblock arXiv:2501.09136v4, 2025.
\newblock \url{https://arxiv.org/abs/2501.09136v4}

\bibitem{arxiv:2411.19443v1}
Tian Yu, Shaolei Zhang, Yang Feng.
\newblock \emph{Auto-RAG: Autonomous Retrieval-Augmented Generation for Large Language Models}.
\newblock arXiv:2411.19443v1, 2024.
\newblock \url{https://arxiv.org/abs/2411.19443v1}

\bibitem{arxiv:2402.01364v2}
Tongtong Wu, Linhao Luo, Yuan-Fang Li, Shirui Pan, Thuy-Trang Vu, Gholamreza Haffari.
\newblock \emph{Continual Learning for Large Language Models A Survey}.
\newblock arXiv:2402.01364v2, 2024.
\newblock \url{https://arxiv.org/abs/2402.01364v2}

\bibitem{arxiv:2603.04647v1}
Xin Chen, Saili Uday Gadgil, Jiarong Qiu.
\newblock \emph{Coordinated Semantic Alignment and Evidence Constraints for Retrieval-Augmented Generation with Large Language Models}.
\newblock arXiv:2603.04647v1, 2026.
\newblock \url{https://arxiv.org/abs/2603.04647v1}

\bibitem{arxiv:2410.05983v1}
Bowen Jin, Jinsung Yoon, Jiawei Han, Sercan O. Arik.
\newblock \emph{Long-Context LLMs Meet RAG: Overcoming Challenges for Long Inputs in RAG}.
\newblock arXiv:2410.05983v1, 2024.
\newblock \url{https://arxiv.org/abs/2410.05983v1}

\bibitem{arxiv:2407.02485v1}
Yue Yu, Wei Ping, Zihan Liu, Boxin Wang, Jiaxuan You, Chao Zhang \emph{et al.}.
\newblock \emph{RankRAG: Unifying Context Ranking with Retrieval-Augmented Generation in LLMs}.
\newblock arXiv:2407.02485v1, 2024.
\newblock \url{https://arxiv.org/abs/2407.02485v1}

\bibitem{arxiv:2501.15915v1}
Weihang Su, Yichen Tang, Qingyao Ai, Junxi Yan, Changyue Wang, Hongning Wang \emph{et al.}.
\newblock \emph{Parametric Retrieval Augmented Generation}.
\newblock arXiv:2501.15915v1, 2025.
\newblock \url{https://arxiv.org/abs/2501.15915v1}

\bibitem{arxiv:2308.07107v3}
Yutao Zhu, Huaying Yuan, Shuting Wang, Jiongnan Liu, Wenhan Liu, Chenlong Deng \emph{et al.}.
\newblock \emph{Large Language Models for Information Retrieval A Survey}.
\newblock arXiv:2308.07107v3, 2023.
\newblock \url{https://arxiv.org/abs/2308.07107v3}

\bibitem{arxiv:2411.19463v3}
Shengming Zhao, Yuchen Shao, Yuheng Huang, Jiayang Song, Zhijie Wang, Chengcheng Wan \emph{et al.}.
\newblock \emph{Understanding the Fundamental Design Decisions of Retrieval-Augmented Generation Systems}.
\newblock arXiv:2411.19463v3, 2024.
\newblock \url{https://arxiv.org/abs/2411.19463v3}

\bibitem{arxiv:2507.18910v1}
Agada Joseph Oche, Ademola Glory Folashade, Tirthankar Ghosal, Arpan Biswas.
\newblock \emph{A Systematic Review of Key Retrieval-Augmented Generation (RAG) Systems: Progress, Gaps, and Future Directions}.
\newblock arXiv:2507.18910v1, 2025.
\newblock \url{https://arxiv.org/abs/2507.18910v1}

\bibitem{arxiv:2503.23895v4}
Yuqiao Tan, Shizhu He, Huanxuan Liao, Jun Zhao, Kang Liu.
\newblock \emph{Dynamic Parametric Retrieval Augmented Generation for Test-time Knowledge Enhancement}.
\newblock arXiv:2503.23895v4, 2025.
\newblock \url{https://arxiv.org/abs/2503.23895v4}

\bibitem{arxiv:2501.07391v1}
Siran Li, Linus Stenzel, Carsten Eickhoff, Seyed Ali Bahrainian.
\newblock \emph{Enhancing Retrieval-Augmented Generation: A Study of Best Practices}.
\newblock arXiv:2501.07391v1, 2025.
\newblock \url{https://arxiv.org/abs/2501.07391v1}

\bibitem{arxiv:2410.07590v1}
Songshuo Lu, Hua Wang, Yutian Rong, Zhi Chen, Yaohua Tang.
\newblock \emph{TurboRAG: Accelerating Retrieval-Augmented Generation with Precomputed KV Caches for Chunked Text}.
\newblock arXiv:2410.07590v1, 2024.
\newblock \url{https://arxiv.org/abs/2410.07590v1}

\bibitem{arxiv:2412.15605v2}
Brian J Chan, Chao-Ting Chen, Jui-Hung Cheng, Hen-Hsen Huang.
\newblock \emph{Don't Do RAG: When Cache-Augmented Generation is All You Need for Knowledge Tasks}.
\newblock arXiv:2412.15605v2, 2024.
\newblock \url{https://arxiv.org/abs/2412.15605v2}

\bibitem{arxiv:2501.13958v3}
Qinggang Zhang, Shengyuan Chen, Yuanchen Bei, Zheng Yuan, Huachi Zhou, Zijin Hong \emph{et al.}.
\newblock \emph{A Survey of Graph Retrieval-Augmented Generation for Customized Large Language Models}.
\newblock arXiv:2501.13958v3, 2025.
\newblock \url{https://arxiv.org/abs/2501.13958v3}

\bibitem{arxiv:2502.14902v2}
Boyu Chen, Zirui Guo, Zidan Yang, Yuluo Chen, Junze Chen, Zhenghao Liu \emph{et al.}.
\newblock \emph{PathRAG: Pruning Graph-based Retrieval Augmented Generation with Relational Paths}.
\newblock arXiv:2502.14902v2, 2025.
\newblock \url{https://arxiv.org/abs/2502.14902v2}

\bibitem{arxiv:2503.19878v3}
Nengbo Wang, Xiaotian Han, Jagdip Singh, Jing Ma, Vipin Chaudhary.
\newblock \emph{CausalRAG: Integrating Causal Graphs into Retrieval-Augmented Generation}.
\newblock arXiv:2503.19878v3, 2025.
\newblock \url{https://arxiv.org/abs/2503.19878v3}

\bibitem{arxiv:2506.20963v2}
Fangyuan Zhang, Zhengjun Huang, Yingli Zhou, Qintian Guo, Zhixun Li, Wensheng Luo \emph{et al.}.
\newblock \emph{EraRAG: Efficient and Incremental Retrieval Augmented Generation for Growing Corpora}.
\newblock arXiv:2506.20963v2, 2025.
\newblock \url{https://arxiv.org/abs/2506.20963v2}

\bibitem{arxiv:2410.10293v3}
Xinping Zhao, Yan Zhong, Zetian Sun, Xinshuo Hu, Zhenyu Liu, Dongfang Li \emph{et al.}.
\newblock \emph{FunnelRAG: A Coarse-to-Fine Progressive Retrieval Paradigm for RAG}.
\newblock arXiv:2410.10293v3, 2024.
\newblock \url{https://arxiv.org/abs/2410.10293v3}

\bibitem{arxiv:2410.10594v2}
Shi Yu, Chaoyue Tang, Bokai Xu, Junbo Cui, Junhao Ran, Yukun Yan \emph{et al.}.
\newblock \emph{VisRAG: Vision-based Retrieval-augmented Generation on Multi-modality Documents}.
\newblock arXiv:2410.10594v2, 2024.
\newblock \url{https://arxiv.org/abs/2410.10594v2}

\bibitem{arxiv:2412.05838v1}
Aniruddha Salve, Saba Attar, Mahesh Deshmukh, Sayali Shivpuje, Arnab Mitra Utsab.
\newblock \emph{A Collaborative Multi-Agent Approach to Retrieval-Augmented Generation Across Diverse Data}.
\newblock arXiv:2412.05838v1, 2024.
\newblock \url{https://arxiv.org/abs/2412.05838v1}

\bibitem{arxiv:2510.12668v2}
Minghao Tang, Shiyu Ni, Jingtong Wu, Zengxin Han, Keping Bi.
\newblock \emph{Understanding Parametric Knowledge Injection in Retrieval-Augmented Generation}.
\newblock arXiv:2510.12668v2, 2025.
\newblock \url{https://arxiv.org/abs/2510.12668v2}

\bibitem{arxiv:2508.06105v2}
Shengyuan Chen, Chuang Zhou, Zheng Yuan, Qinggang Zhang, Zeyang Cui, Hao Chen \emph{et al.}.
\newblock \emph{You Don't Need Pre-built Graphs for RAG: Retrieval Augmented Generation with Adaptive Reasoning Structures}.
\newblock arXiv:2508.06105v2, 2025.
\newblock \url{https://arxiv.org/abs/2508.06105v2}

\bibitem{arxiv:2511.05385v2}
Chao Zhang, Yuhao Wang, Derong Xu, Haoxin Zhang, Yuanjie Lyu, Yuhao Chen \emph{et al.}.
\newblock \emph{TeaRAG: A Token-Efficient Agentic Retrieval-Augmented Generation Framework}.
\newblock arXiv:2511.05385v2, 2025.
\newblock \url{https://arxiv.org/abs/2511.05385v2}

\bibitem{arxiv:2604.12610v1}
Xudong Wang, Chaoning Zhang, Qigan Sun, Zhenzhen Huang, Chang Lu, Sheng Zheng \emph{et al.}.
\newblock \emph{Transforming External Knowledge into Triplets for Enhanced Retrieval in RAG of LLMs}.
\newblock arXiv:2604.12610v1, 2026.
\newblock \url{https://arxiv.org/abs/2604.12610v1}

\bibitem{arxiv:2501.18365v1}
Yiteng Tu, Weihang Su, Yujia Zhou, Yiqun Liu, Qingyao Ai.
\newblock \emph{RbFT: Robust Fine-tuning for Retrieval-Augmented Generation against Retrieval Defects}.
\newblock arXiv:2501.18365v1, 2025.
\newblock \url{https://arxiv.org/abs/2501.18365v1}

\bibitem{arxiv:2410.13339v2}
Ingeol Baek, Hwan Chang, Byeongjeong Kim, Jimin Lee, Hwanhee Lee.
\newblock \emph{Probing-RAG: Self-Probing to Guide Language Models in Selective Document Retrieval}.
\newblock arXiv:2410.13339v2, 2024.
\newblock \url{https://arxiv.org/abs/2410.13339v2}

\bibitem{arxiv:2512.20144v2}
Yuxin Wang, Shicheng Fang, Bo Wang, Qi Luo, Xuanjing Huang, Yining Zheng \emph{et al.}.
\newblock \emph{Multi-hop Reasoning via Early Knowledge Alignment}.
\newblock arXiv:2512.20144v2, 2025.
\newblock \url{https://arxiv.org/abs/2512.20144v2}

\bibitem{arxiv:2512.03413v1}
Shu Wang, Yingli Zhou, Yixiang Fang.
\newblock \emph{BookRAG: A Hierarchical Structure-aware Index-based Approach for Retrieval-Augmented Generation on Complex Documents}.
\newblock arXiv:2512.03413v1, 2025.
\newblock \url{https://arxiv.org/abs/2512.03413v1}

\bibitem{arxiv:2412.05547v2}
Weijie Chen, Ting Bai, Jinbo Su, Jian Luan, Wei Liu, Chuan Shi.
\newblock \emph{KG-Retriever: Efficient Knowledge Indexing for Retrieval-Augmented Large Language Models}.
\newblock arXiv:2412.05547v2, 2024.
\newblock \url{https://arxiv.org/abs/2412.05547v2}

\bibitem{arxiv:2501.09957v2}
Zengyi Gao, Yukun Cao, Hairu Wang, Ao Ke, Yuan Feng, Xike Xie \emph{et al.}.
\newblock \emph{FRAG: A Flexible Modular Framework for Retrieval-Augmented Generation based on Knowledge Graphs}.
\newblock arXiv:2501.09957v2, 2025.
\newblock \url{https://arxiv.org/abs/2501.09957v2}

\bibitem{arxiv:2411.11895v1}
Sonal Prabhune, Donald J. Berndt.
\newblock \emph{Deploying Large Language Models With Retrieval Augmented Generation}.
\newblock arXiv:2411.11895v1, 2024.
\newblock \url{https://arxiv.org/abs/2411.11895v1}

\bibitem{arxiv:2505.23794v2}
Yuan Li, Qi Luo, Xiaonan Li, Bufan Li, Qinyuan Cheng, Bo Wang \emph{et al.}.
\newblock \emph{R3-RAG: Learning Step-by-Step Reasoning and Retrieval for LLMs via Reinforcement Learning}.
\newblock arXiv:2505.23794v2, 2025.
\newblock \url{https://arxiv.org/abs/2505.23794v2}

\bibitem{arxiv:2502.19280v2}
Akash Dhasade, Rachid Guerraoui, Anne-Marie Kermarrec, Diana Petrescu, Rafael Pires, Mathis Randl \emph{et al.}.
\newblock \emph{Efficient Federated Search for Retrieval-Augmented Generation using Lightweight Routing}.
\newblock arXiv:2502.19280v2, 2025.
\newblock \url{https://arxiv.org/abs/2502.19280v2}

\bibitem{arxiv:2502.11371v3}
Haoyu Han, Li Ma, Yu Wang, Harry Shomer, Yongjia Lei, Zhisheng Qi \emph{et al.}.
\newblock \emph{RAG vs. GraphRAG: A Systematic Evaluation and Key Insights}.
\newblock arXiv:2502.11371v3, 2025.
\newblock \url{https://arxiv.org/abs/2502.11371v3}

\bibitem{arxiv:2405.06211v3}
Wenqi Fan, Yujuan Ding, Liangbo Ning, Shijie Wang, Hengyun Li, Dawei Yin \emph{et al.}.
\newblock \emph{A Survey on RAG Meeting LLMs: Towards Retrieval-Augmented Large Language Models}.
\newblock arXiv:2405.06211v3, 2024.
\newblock \url{https://arxiv.org/abs/2405.06211v3}

\bibitem{arxiv:2310.11511v1}
Akari Asai, Zeqiu Wu, Yizhong Wang, Avirup Sil, Hannaneh Hajishirzi.
\newblock \emph{Self-RAG Learning to Retrieve, Generate, and Critique through Self-Reflection}.
\newblock arXiv:2310.11511v1, 2023.
\newblock \url{https://arxiv.org/abs/2310.11511v1}

\bibitem{arxiv:2508.06401v3}
Andrew Brown, Muhammad Roman, Barry Devereux.
\newblock \emph{A Systematic Literature Review of Retrieval-Augmented Generation: Techniques, Metrics, and Challenges}.
\newblock arXiv:2508.06401v3, 2025.
\newblock \url{https://arxiv.org/abs/2508.06401v3}

\bibitem{arxiv:2604.00715v1}
Karan Singh, Michael Yu, Varun Gangal, Zhuofu Tao, Sachin Kumar, Emmy Liu \emph{et al.}.
\newblock \emph{To Memorize or to Retrieve: Scaling Laws for RAG-Considerate Pretraining}.
\newblock arXiv:2604.00715v1, 2026.
\newblock \url{https://arxiv.org/abs/2604.00715v1}

\bibitem{arxiv:2401.14887v3}
Florin Cuconasu, Giovanni Trappolini, Federico Siciliano, Simone Filice, Cesare Campagnano, Yoelle Maarek \emph{et al.}.
\newblock \emph{The Power of Noise Redefining Retrieval for RAG Systems}.
\newblock arXiv:2401.14887v3, 2024.
\newblock \url{https://arxiv.org/abs/2401.14887v3}

\bibitem{arxiv:2312.05934v3}
Oded Ovadia, Menachem Brief, Moshik Mishaeli, Oren Elisha.
\newblock \emph{Fine-Tuning or Retrieval Comparing Knowledge Injection in LLMs}.
\newblock arXiv:2312.05934v3, 2023.
\newblock \url{https://arxiv.org/abs/2312.05934v3}

\bibitem{arxiv:2410.05779v3}
Zirui Guo, Lianghao Xia, Yanhua Yu, Tu Ao, Chao Huang.
\newblock \emph{LightRAG: Simple and Fast Retrieval-Augmented Generation}.
\newblock arXiv:2410.05779v3, 2024.
\newblock \url{https://arxiv.org/abs/2410.05779v3}

\bibitem{arxiv:2604.10426v1}
Cheng-Yen Li, Xuanjun Chen, Claire Lin, Wei-Yu Chen, Wenhua Nie, Hung-Yi Lee \emph{et al.}.
\newblock \emph{CodaRAG: Connecting the Dots with Associativity Inspired by Complementary Learning}.
\newblock arXiv:2604.10426v1, 2026.
\newblock \url{https://arxiv.org/abs/2604.10426v1}

\bibitem{arxiv:2412.15529v4}
Qili Zhang, Qianren Mao, Yangyifei Luo, Yashuo Luo, Hanwen Hao, Zhilong Cao \emph{et al.}.
\newblock \emph{XRAG: eXamining the Core -- Benchmarking Foundational Components in Advanced Retrieval-Augmented Generation}.
\newblock arXiv:2412.15529v4, 2024.
\newblock \url{https://arxiv.org/abs/2412.15529v4}

\bibitem{arxiv:2404.10981v1}
Yizheng Huang, Jimmy Huang.
\newblock \emph{A Survey on Retrieval-Augmented Text Generation for Large Language Models}.
\newblock arXiv:2404.10981v1, 2024.
\newblock \url{https://arxiv.org/abs/2404.10981v1}

\bibitem{arxiv:2504.19754v1}
Carlo Merola, Jaspinder Singh.
\newblock \emph{Reconstructing Context: Evaluating Advanced Chunking Strategies for Retrieval-Augmented Generation}.
\newblock arXiv:2504.19754v1, 2025.
\newblock \url{https://arxiv.org/abs/2504.19754v1}

\bibitem{arxiv:2403.14403v2}
Soyeong Jeong, Jinheon Baek, Sukmin Cho, Sung Ju Hwang, Jong C. Park.
\newblock \emph{Adaptive-RAG Learning to Adapt Retrieval-Augmented Large Language Models through Question Complexity}.
\newblock arXiv:2403.14403v2, 2024.
\newblock \url{https://arxiv.org/abs/2403.14403v2}

\bibitem{arxiv:2412.01572v4}
Xiaqiang Tang, Qiang Gao, Jian Li, Nan Du, Qi Li, Sihong Xie.
\newblock \emph{MBA-RAG: a Bandit Approach for Adaptive Retrieval-Augmented Generation through Question Complexity}.
\newblock arXiv:2412.01572v4, 2024.
\newblock \url{https://arxiv.org/abs/2412.01572v4}

\bibitem{arxiv:2604.15771v3}
Kai Wei, Raymond Li, Xi Zhu, Zhaoqian Xue, Jiaojiao Han, Jingcheng Niu \emph{et al.}.
\newblock \emph{Skill-RAG: Failure-State-Aware Retrieval Augmentation via Hidden-State Probing and Skill Routing}.
\newblock arXiv:2604.15771v3, 2026.
\newblock \url{https://arxiv.org/abs/2604.15771v3}

\bibitem{arxiv:2505.14146v2}
Pengcheng Jiang, Xueqiang Xu, Jiacheng Lin, Jinfeng Xiao, Zifeng Wang, Jimeng Sun \emph{et al.}.
\newblock \emph{s3: You Don't Need That Much Data to Train a Search Agent via RL}.
\newblock arXiv:2505.14146v2, 2025.
\newblock \url{https://arxiv.org/abs/2505.14146v2}

\bibitem{arxiv:2505.02811v2}
Diji Yang, Linda Zeng, Jinmeng Rao, Yi Zhang.
\newblock \emph{Knowing You Don't Know: Learning When to Continue Search in Multi-round RAG through Self-Practicing}.
\newblock arXiv:2505.02811v2, 2025.
\newblock \url{https://arxiv.org/abs/2505.02811v2}

\bibitem{arxiv:2602.03442v1}
Mingxuan Du, Benfeng Xu, Chiwei Zhu, Shaohan Wang, Pengyu Wang, Xiaorui Wang \emph{et al.}.
\newblock \emph{A-RAG: Scaling Agentic Retrieval-Augmented Generation via Hierarchical Retrieval Interfaces}.
\newblock arXiv:2602.03442v1, 2026.
\newblock \url{https://arxiv.org/abs/2602.03442v1}

\bibitem{arxiv:2410.07176v2}
Fei Wang, Xingchen Wan, Ruoxi Sun, Jiefeng Chen, Sercan \"{O}. Ar\i{}k.
\newblock \emph{Astute RAG: Overcoming Imperfect Retrieval Augmentation and Knowledge Conflicts for Large Language Models}.
\newblock arXiv:2410.07176v2, 2024.
\newblock \url{https://arxiv.org/abs/2410.07176v2}

\bibitem{arxiv:2410.19572v5}
Ishneet Sukhvinder Singh, Ritvik Aggarwal, Ibrahim Allahverdiyev, Muhammad Taha, Aslihan Akalin, Kevin Zhu \emph{et al.}.
\newblock \emph{ChunkRAG: Novel LLM-Chunk Filtering Method for RAG Systems}.
\newblock arXiv:2410.19572v5, 2024.
\newblock \url{https://arxiv.org/abs/2410.19572v5}

\bibitem{arxiv:2306.05212v1}
Jiongnan Liu, Jiajie Jin, Zihan Wang, Jiehan Cheng, Zhicheng Dou, Ji-Rong Wen.
\newblock \emph{RETA-LLM A Retrieval-Augmented Large Language Model Toolkit}.
\newblock arXiv:2306.05212v1, 2023.
\newblock \url{https://arxiv.org/abs/2306.05212v1}

\bibitem{arxiv:2401.06954v2}
Zixuan Ke, Weize Kong, Cheng Li, Mingyang Zhang, Qiaozhu Mei, Michael Bendersky.
\newblock \emph{Bridging the Preference Gap between Retrievers and LLMs}.
\newblock arXiv:2401.06954v2, 2024.
\newblock \url{https://arxiv.org/abs/2401.06954v2}

\bibitem{arxiv:2507.07030v1}
Fengran Mo, Yifan Gao, Chuan Meng, Xin Liu, Zhuofeng Wu, Kelong Mao \emph{et al.}.
\newblock \emph{UniConv: Unifying Retrieval and Response Generation for Large Language Models in Conversations}.
\newblock arXiv:2507.07030v1, 2025.
\newblock \url{https://arxiv.org/abs/2507.07030v1}

\bibitem{arxiv:2605.08333v1}
Pengzhou Chen, Tao Chen.
\newblock \emph{CDS4RAG: Cyclic Dual-Sequential Hyperparameter Optimization for RAG}.
\newblock arXiv:2605.08333v1, 2026.
\newblock \url{https://arxiv.org/abs/2605.08333v1}

\bibitem{arxiv:2504.15909v2}
Yunfan Gao, Yun Xiong, Yijie Zhong, Yuxi Bi, Ming Xue, Haofen Wang.
\newblock \emph{Synergizing RAG and Reasoning: A Systematic Review}.
\newblock arXiv:2504.15909v2, 2025.
\newblock \url{https://arxiv.org/abs/2504.15909v2}

\bibitem{arxiv:2503.04338v2}
Yingli Zhou, Yaodong Su, Youran Sun, Shu Wang, Taotao Wang, Runyuan He \emph{et al.}.
\newblock \emph{In-depth Analysis of Graph-based RAG in a Unified Framework}.
\newblock arXiv:2503.04338v2, 2025.
\newblock \url{https://arxiv.org/abs/2503.04338v2}

\bibitem{arxiv:2601.08620v2}
Ant\'{o}nio Loison, Quentin Mac\'{e}, Antoine Edy, Victor Xing, Tom Balough, Gabriel Moreira \emph{et al.}.
\newblock \emph{ViDoRe V3: A Comprehensive Evaluation of Retrieval Augmented Generation in Complex Real-World Scenarios}.
\newblock arXiv:2601.08620v2, 2026.
\newblock \url{https://arxiv.org/abs/2601.08620v2}

\bibitem{arxiv:2502.15734v1}
Shubham Agarwal, Sai Sundaresan, Subrata Mitra, Debabrata Mahapatra, Archit Gupta, Rounak Sharma \emph{et al.}.
\newblock \emph{Cache-Craft: Managing Chunk-Caches for Efficient Retrieval-Augmented Generation}.
\newblock arXiv:2502.15734v1, 2025.
\newblock \url{https://arxiv.org/abs/2502.15734v1}

\bibitem{arxiv:2309.15217v1}
Shahul Es, Jithin James, Luis Espinosa-Anke, Steven Schockaert.
\newblock \emph{RAGAS Automated Evaluation of Retrieval Augmented Generation}.
\newblock arXiv:2309.15217v1, 2023.
\newblock \url{https://arxiv.org/abs/2309.15217v1}

\bibitem{arxiv:2603.01783v1}
Yifan Wang, Mingxuan Jiang, Zhihao Sun, Yixin Cao, Yicun Liu, Keyang Chen \emph{et al.}.
\newblock \emph{GAM-RAG: Gain-Adaptive Memory for Evolving Retrieval in Retrieval-Augmented Generation}.
\newblock arXiv:2603.01783v1, 2026.
\newblock \url{https://arxiv.org/abs/2603.01783v1}

\bibitem{arxiv:2404.12457v2}
Chao Jin, Zili Zhang, Xuanlin Jiang, Fangyue Liu, Xin Liu, Xuanzhe Liu \emph{et al.}.
\newblock \emph{RAGCache Efficient Knowledge Caching for Retrieval-Augmented Generation}.
\newblock arXiv:2404.12457v2, 2024.
\newblock \url{https://arxiv.org/abs/2404.12457v2}

\bibitem{arxiv:2602.02579v3}
Shihao Wang, Jiahao Chen, Yanqi Pan, Hao Huang, Yichen Hao, Xiangyu Zou \emph{et al.}.
\newblock \emph{ProphetKV: User-Query-Driven Selective Recomputation for Efficient KV Cache Reuse in Retrieval-Augmented Generation}.
\newblock arXiv:2602.02579v3, 2026.
\newblock \url{https://arxiv.org/abs/2602.02579v3}

\bibitem{arxiv:2507.09477v2}
Yangning Li, Weizhi Zhang, Yuyao Yang, Wei-Chieh Huang, Yaozu Wu, Junyu Luo \emph{et al.}.
\newblock \emph{Towards Agentic RAG with Deep Reasoning: A Survey of RAG-Reasoning Systems in LLMs}.
\newblock arXiv:2507.09477v2, 2025.
\newblock \url{https://arxiv.org/abs/2507.09477v2}

\bibitem{arxiv:2506.06704v1}
Weihang Su, Qingyao Ai, Jingtao Zhan, Qian Dong, Yiqun Liu.
\newblock \emph{Dynamic and Parametric Retrieval-Augmented Generation}.
\newblock arXiv:2506.06704v1, 2025.
\newblock \url{https://arxiv.org/abs/2506.06704v1}

\bibitem{arxiv:2411.14572v1}
Shenglai Zeng, Jiankun Zhang, Bingheng Li, Yuping Lin, Tianqi Zheng, Dante Everaert \emph{et al.}.
\newblock \emph{Towards Knowledge Checking in Retrieval-augmented Generation: A Representation Perspective}.
\newblock arXiv:2411.14572v1, 2024.
\newblock \url{https://arxiv.org/abs/2411.14572v1}

\bibitem{arxiv:2505.17118v2}
Xinbang Dai, Huikang Hu, Yuncheng Hua, Jiaqi Li, Yongrui Chen, Rihui Jin \emph{et al.}.
\newblock \emph{After Retrieval, Before Generation: Enhancing the Trustworthiness of Large Language Models in Retrieval-Augmented Generation}.
\newblock arXiv:2505.17118v2, 2025.
\newblock \url{https://arxiv.org/abs/2505.17118v2}

\bibitem{arxiv:2305.15294v2}
Zhihong Shao, Yeyun Gong, Yelong Shen, Minlie Huang, Nan Duan, Weizhu Chen.
\newblock \emph{Enhancing Retrieval-Augmented Large Language Models with Iterative Retrieval-Generation Synergy}.
\newblock arXiv:2305.15294v2, 2023.
\newblock \url{https://arxiv.org/abs/2305.15294v2}

\bibitem{arxiv:2505.17005v1}
Huatong Song, Jinhao Jiang, Wenqing Tian, Zhipeng Chen, Yuhuan Wu, Jiahao Zhao \emph{et al.}.
\newblock \emph{R1-Searcher++: Incentivizing the Dynamic Knowledge Acquisition of LLMs via Reinforcement Learning}.
\newblock arXiv:2505.17005v1, 2025.
\newblock \url{https://arxiv.org/abs/2505.17005v1}

\bibitem{arxiv:2502.01142v2}
Xinyan Guan, Jiali Zeng, Fandong Meng, Chunlei Xin, Yaojie Lu, Hongyu Lin \emph{et al.}.
\newblock \emph{DeepRAG: Thinking to Retrieve Step by Step for Large Language Models}.
\newblock arXiv:2502.01142v2, 2025.
\newblock \url{https://arxiv.org/abs/2502.01142v2}

\bibitem{arxiv:2410.05162v1}
Mehrdad Farahani, Richard Johansson.
\newblock \emph{Deciphering the Interplay of Parametric and Non-parametric Memory in Retrieval-augmented Language Models}.
\newblock arXiv:2410.05162v1, 2024.
\newblock \url{https://arxiv.org/abs/2410.05162v1}

\bibitem{arxiv:2503.14649v2}
Wenqi Jiang, Suvinay Subramanian, Cat Graves, Gustavo Alonso, Amir Yazdanbakhsh, Vidushi Dadu.
\newblock \emph{RAGO: Systematic Performance Optimization for Retrieval-Augmented Generation Serving}.
\newblock arXiv:2503.14649v2, 2025.
\newblock \url{https://arxiv.org/abs/2503.14649v2}

\bibitem{arxiv:2504.08893v1}
Jasper Linders, Jakub M. Tomczak.
\newblock \emph{Knowledge Graph-extended Retrieval Augmented Generation for Question Answering}.
\newblock arXiv:2504.08893v1, 2025.
\newblock \url{https://arxiv.org/abs/2504.08893v1}

\bibitem{arxiv:2506.10408v1}
Jintao Liang, Gang Su, Huifeng Lin, You Wu, Rui Zhao, Ziyue Li.
\newblock \emph{Reasoning RAG via System 1 or System 2: A Survey on Reasoning Agentic Retrieval-Augmented Generation for Industry Challenges}.
\newblock arXiv:2506.10408v1, 2025.
\newblock \url{https://arxiv.org/abs/2506.10408v1}

\bibitem{arxiv:2501.00332v1}
Chia-Yuan Chang, Zhimeng Jiang, Vineeth Rakesh, Menghai Pan, Chin-Chia Michael Yeh, Guanchu Wang \emph{et al.}.
\newblock \emph{MAIN-RAG: Multi-Agent Filtering Retrieval-Augmented Generation}.
\newblock arXiv:2501.00332v1, 2025.
\newblock \url{https://arxiv.org/abs/2501.00332v1}

\bibitem{arxiv:2509.03934v1}
Yuqing Huang, Rongyang Zhang, Qimeng Wang, Chengqiang Lu, Yan Gao, Yi Wu \emph{et al.}.
\newblock \emph{SelfAug: Mitigating Catastrophic Forgetting in Retrieval-Augmented Generation via Distribution Self-Alignment}.
\newblock arXiv:2509.03934v1, 2025.
\newblock \url{https://arxiv.org/abs/2509.03934v1}

\bibitem{arxiv:2510.13910v2}
Jingru Lin, Chen Zhang, Stephen Y. Liu, Haizhou Li.
\newblock \emph{RAGCap-Bench: Benchmarking Capabilities of LLMs in Agentic Retrieval Augmented Generation Systems}.
\newblock arXiv:2510.13910v2, 2025.
\newblock \url{https://arxiv.org/abs/2510.13910v2}

\bibitem{arxiv:2601.06842v1}
Hua Ye, Siyuan Chen, Ziqi Zhong, Canran Xiao, Haoliang Zhang, Yuhan Wu \emph{et al.}.
\newblock \emph{Seeing through the Conflict: Transparent Knowledge Conflict Handling in Retrieval-Augmented Generation}.
\newblock arXiv:2601.06842v1, 2026.
\newblock \url{https://arxiv.org/abs/2601.06842v1}

\bibitem{arxiv:2506.20051v2}
Jia-Huei Ju, Suzan Verberne, Maarten de Rijke, Andrew Yates.
\newblock \emph{Controlled Retrieval-augmented Context Evaluation for Long-form RAG}.
\newblock arXiv:2506.20051v2, 2025.
\newblock \url{https://arxiv.org/abs/2506.20051v2}

\bibitem{arxiv:2309.01431v2}
Jiawei Chen, Hongyu Lin, Xianpei Han, Le Sun.
\newblock \emph{Benchmarking Large Language Models in Retrieval-Augmented Generation}.
\newblock arXiv:2309.01431v2, 2023.
\newblock \url{https://arxiv.org/abs/2309.01431v2}

\bibitem{arxiv:2410.12248v1}
Jintao Liu, Ruixue Ding, Linhao Zhang, Pengjun Xie, Fie Huang.
\newblock \emph{CoFE-RAG: A Comprehensive Full-chain Evaluation Framework for Retrieval-Augmented Generation with Enhanced Data Diversity}.
\newblock arXiv:2410.12248v1, 2024.
\newblock \url{https://arxiv.org/abs/2410.12248v1}

\bibitem{arxiv:2412.11854v1}
Michael Shen, Muhammad Umar, Kiwan Maeng, G. Edward Suh, Udit Gupta.
\newblock \emph{Towards Understanding Systems Trade-offs in Retrieval-Augmented Generation Model Inference}.
\newblock arXiv:2412.11854v1, 2024.
\newblock \url{https://arxiv.org/abs/2412.11854v1}

\bibitem{arxiv:2605.15156v2}
Ryan Wei Heng Quek, Sanghyuk Lee, Alfred Wei Lun Leong, Arun Verma, Alok Prakash, Nancy F. Chen \emph{et al.}.
\newblock \emph{MeMo: Memory as a Model}.
\newblock arXiv:2605.15156v2, 2026.
\newblock \url{https://arxiv.org/abs/2605.15156v2}

\bibitem{arxiv:2503.18016v1}
Xu Zheng, Ziqiao Weng, Yuanhuiyi Lyu, Lutao Jiang, Haiwei Xue, Bin Ren \emph{et al.}.
\newblock \emph{Retrieval Augmented Generation and Understanding in Vision: A Survey and New Outlook}.
\newblock arXiv:2503.18016v1, 2025.
\newblock \url{https://arxiv.org/abs/2503.18016v1}

\bibitem{arxiv:2503.10150v3}
Haoyu Huang, Yongfeng Huang, Junjie Yang, Zhenyu Pan, Yongqiang Chen, Kaili Ma \emph{et al.}.
\newblock \emph{Retrieval-Augmented Generation with Hierarchical Knowledge}.
\newblock arXiv:2503.10150v3, 2025.
\newblock \url{https://arxiv.org/abs/2503.10150v3}

\bibitem{arxiv:2502.06864v1}
Xiangrong Zhu, Yuexiang Xie, Yi Liu, Yaliang Li, Wei Hu.
\newblock \emph{Knowledge Graph-Guided Retrieval Augmented Generation}.
\newblock arXiv:2502.06864v1, 2025.
\newblock \url{https://arxiv.org/abs/2502.06864v1}

\bibitem{arxiv:2504.10074v3}
Zihan Ling, Zhiyao Guo, Yixuan Huang, Yi An, Shuai Xiao, Jinsong Lan \emph{et al.}.
\newblock \emph{MMKB-RAG: A Multi-Modal Knowledge-Based Retrieval-Augmented Generation Framework}.
\newblock arXiv:2504.10074v3, 2025.
\newblock \url{https://arxiv.org/abs/2504.10074v3}

\bibitem{arxiv:2502.06872v1}
Bo Ni, Zheyuan Liu, Leyao Wang, Yongjia Lei, Yuying Zhao, Xueqi Cheng \emph{et al.}.
\newblock \emph{Towards Trustworthy Retrieval Augmented Generation for Large Language Models: A Survey}.
\newblock arXiv:2502.06872v1, 2025.
\newblock \url{https://arxiv.org/abs/2502.06872v1}

\bibitem{arxiv:2502.14759v1}
Juraj Vladika, Florian Matthes.
\newblock \emph{On the Influence of Context Size and Model Choice in Retrieval-Augmented Generation Systems}.
\newblock arXiv:2502.14759v1, 2025.
\newblock \url{https://arxiv.org/abs/2502.14759v1}

\bibitem{arxiv:2502.01113v3}
Linhao Luo, Zicheng Zhao, Gholamreza Haffari, Dinh Phung, Chen Gong, Shirui Pan.
\newblock \emph{GFM-RAG: Graph Foundation Model for Retrieval Augmented Generation}.
\newblock arXiv:2502.01113v3, 2025.
\newblock \url{https://arxiv.org/abs/2502.01113v3}

\bibitem{arxiv:2401.15884v2}
Shi-Qi Yan, Jia-Chen Gu, Yun Zhu, Zhen-Hua Ling.
\newblock \emph{Corrective Retrieval Augmented Generation}.
\newblock arXiv:2401.15884v2, 2024.
\newblock \url{https://arxiv.org/abs/2401.15884v2}

\bibitem{arxiv:2505.03275v1}
Tiantian Gan, Qiyao Sun.
\newblock \emph{RAG-MCP: Mitigating Prompt Bloat in LLM Tool Selection via Retrieval-Augmented Generation}.
\newblock arXiv:2505.03275v1, 2025.
\newblock \url{https://arxiv.org/abs/2505.03275v1}

\bibitem{arxiv:2401.05856v1}
Scott Barnett, Stefanus Kurniawan, Srikanth Thudumu, Zach Brannelly, Mohamed Abdelrazek.
\newblock \emph{Seven Failure Points When Engineering a Retrieval Augmented Generation System}.
\newblock arXiv:2401.05856v1, 2024.
\newblock \url{https://arxiv.org/abs/2401.05856v1}

\bibitem{arxiv:2510.02657v3}
Jingjie Ning, Yibo Kong, Yunfan Long, Jamie Callan.
\newblock \emph{Less LLM, More Documents: Searching for Improved RAG}.
\newblock arXiv:2510.02657v3, 2025.
\newblock \url{https://arxiv.org/abs/2510.02657v3}

\bibitem{arxiv:2208.07652v1}
Jiangui Chen, Ruqing Zhang, Jiafeng Guo, Yiqun Liu, Yixing Fan, Xueqi Cheng.
\newblock \emph{CorpusBrain Pre-train a Generative Retrieval Model for Knowledge-Intensive Language Tasks}.
\newblock arXiv:2208.07652v1, 2022.
\newblock \url{https://arxiv.org/abs/2208.07652v1}

\bibitem{arxiv:2502.00592v2}
Yu Wang, Dmitry Krotov, Yuanzhe Hu, Yifan Gao, Wangchunshu Zhou, Julian McAuley \emph{et al.}.
\newblock \emph{M+: Extending MemoryLLM with Scalable Long-Term Memory}.
\newblock arXiv:2502.00592v2, 2025.
\newblock \url{https://arxiv.org/abs/2502.00592v2}

\bibitem{arxiv:2412.18431v2}
Zhili Shen, Chenxin Diao, Pavlos Vougiouklis, Pascual Merita, Shriram Piramanayagam, Enting Chen \emph{et al.}.
\newblock \emph{GeAR: Graph-enhanced Agent for Retrieval-augmented Generation}.
\newblock arXiv:2412.18431v2, 2024.
\newblock \url{https://arxiv.org/abs/2412.18431v2}

\bibitem{arxiv:2511.05549v2}
Yubo Wang, Haoyang Li, Fei Teng, Lei Chen.
\newblock \emph{AGRAG: Advanced Graph-based Retrieval-Augmented Generation for LLMs}.
\newblock arXiv:2511.05549v2, 2025.
\newblock \url{https://arxiv.org/abs/2511.05549v2}

\bibitem{arxiv:2502.14727v1}
Yifu Chen, Shengpeng Ji, Haoxiao Wang, Ziqing Wang, Siyu Chen, Jinzheng He \emph{et al.}.
\newblock \emph{WavRAG: Audio-Integrated Retrieval Augmented Generation for Spoken Dialogue Models}.
\newblock arXiv:2502.14727v1, 2025.
\newblock \url{https://arxiv.org/abs/2502.14727v1}

\bibitem{arxiv:2409.19401v1}
Zheng Wang, Zhongyang Li, Zeren Jiang, Dandan Tu, Wei Shi.
\newblock \emph{Crafting Personalized Agents through Retrieval-Augmented Generation on Editable Memory Graphs}.
\newblock arXiv:2409.19401v1, 2024.
\newblock \url{https://arxiv.org/abs/2409.19401v1}

\bibitem{arxiv:2503.10677v3}
Mingyue Cheng, Yucong Luo, Jie Ouyang, Qi Liu, Huijie Liu, Li Li \emph{et al.}.
\newblock \emph{A Survey on Knowledge-Oriented Retrieval-Augmented Generation}.
\newblock arXiv:2503.10677v3, 2025.
\newblock \url{https://arxiv.org/abs/2503.10677v3}

\bibitem{arxiv:2506.21931v2}
Reza Yousefi Maragheh, Pratheek Vadla, Priyank Gupta, Kai Zhao, Aysenur Inan, Kehui Yao \emph{et al.}.
\newblock \emph{ARAG: Agentic Retrieval Augmented Generation for Personalized Recommendation}.
\newblock arXiv:2506.21931v2, 2025.
\newblock \url{https://arxiv.org/abs/2506.21931v2}

\bibitem{arxiv:2602.03578v1}
Su Dong, Qinggang Zhang, Yilin Xiao, Shengyuan Chen, Chuang Zhou, Xiao Huang.
\newblock \emph{Use Graph When It Needs: Efficiently and Adaptively Integrating Retrieval-Augmented Generation with Graphs}.
\newblock arXiv:2602.03578v1, 2026.
\newblock \url{https://arxiv.org/abs/2602.03578v1}

\bibitem{arxiv:2510.13975v2}
Kin Kwan Leung, Mouloud Belbahri, Yi Sui, Alex Labach, Xueying Zhang, Stephen Anthony Rose \emph{et al.}.
\newblock \emph{Classifying and Addressing the Diversity of Errors in Retrieval-Augmented Generation Systems}.
\newblock arXiv:2510.13975v2, 2025.
\newblock \url{https://arxiv.org/abs/2510.13975v2}

\bibitem{arxiv:2401.15391v1}
Yixuan Tang, Yi Yang.
\newblock \emph{MultiHop-RAG Benchmarking Retrieval-Augmented Generation for Multi-Hop Queries}.
\newblock arXiv:2401.15391v1, 2024.
\newblock \url{https://arxiv.org/abs/2401.15391v1}

\bibitem{arxiv:2502.15543v4}
Pengcheng Huang, Zhenghao Liu, Yukun Yan, Haiyan Zhao, Xiaoyuan Yi, Hao Chen \emph{et al.}.
\newblock \emph{ParamMute: Suppressing Knowledge-Critical FFNs for Faithful Retrieval-Augmented Generation}.
\newblock arXiv:2502.15543v4, 2025.
\newblock \url{https://arxiv.org/abs/2502.15543v4}

\bibitem{arxiv:2506.08938v2}
Qinggang Zhang, Zhishang Xiang, Yilin Xiao, Le Wang, Junhui Li, Xinrun Wang \emph{et al.}.
\newblock \emph{FaithfulRAG: Fact-Level Conflict Modeling for Context-Faithful Retrieval-Augmented Generation}.
\newblock arXiv:2506.08938v2, 2025.
\newblock \url{https://arxiv.org/abs/2506.08938v2}

\bibitem{arxiv:2511.10375v1}
Shuyi Liu, Yuming Shang, Xi Zhang.
\newblock \emph{TruthfulRAG: Resolving Factual-level Conflicts in Retrieval-Augmented Generation with Knowledge Graphs}.
\newblock arXiv:2511.10375v1, 2025.
\newblock \url{https://arxiv.org/abs/2511.10375v1}

\bibitem{arxiv:2410.03439v3}
Renxi Wang, Xudong Han, Lei Ji, Shu Wang, Timothy Baldwin, Haonan Li.
\newblock \emph{ToolGen: Unified Tool Retrieval and Calling via Generation}.
\newblock arXiv:2410.03439v3, 2024.
\newblock \url{https://arxiv.org/abs/2410.03439v3}

\bibitem{arxiv:2502.18635v2}
Matthew Barker, Andrew Bell, Evan Thomas, James Carr, Thomas Andrews, Umang Bhatt.
\newblock \emph{Faster, Cheaper, Better: Multi-Objective Hyperparameter Optimization for LLM and RAG Systems}.
\newblock arXiv:2502.18635v2, 2025.
\newblock \url{https://arxiv.org/abs/2502.18635v2}

\bibitem{arxiv:2410.01782v1}
Shayekh Bin Islam, Md Asib Rahman, K S M Tozammel Hossain, Enamul Hoque, Shafiq Joty, Md Rizwan Parvez.
\newblock \emph{Open-RAG: Enhanced Retrieval-Augmented Reasoning with Open-Source Large Language Models}.
\newblock arXiv:2410.01782v1, 2024.
\newblock \url{https://arxiv.org/abs/2410.01782v1}

\bibitem{arxiv:2110.07752v2}
Ashwin Paranjape, Omar Khattab, Christopher Potts, Matei Zaharia, Christopher D. Manning.
\newblock \emph{Hindsight Posterior-guided training of retrievers for improved open-ended generation}.
\newblock arXiv:2110.07752v2, 2021.
\newblock \url{https://arxiv.org/abs/2110.07752v2}

\bibitem{arxiv:2509.23519v2}
Zeyu Shen, Basileal Imana, Tong Wu, Chong Xiang, Prateek Mittal, Aleksandra Korolova.
\newblock \emph{ReliabilityRAG: Effective and Provably Robust Defense for RAG-based Web-Search}.
\newblock arXiv:2509.23519v2, 2025.
\newblock \url{https://arxiv.org/abs/2509.23519v2}

\bibitem{arxiv:2502.14802v2}
Bernal Jim\'{e}nez Guti\'{e}rrez, Yiheng Shu, Weijian Qi, Sizhe Zhou, Yu Su.
\newblock \emph{From RAG to Memory: Non-Parametric Continual Learning for Large Language Models}.
\newblock arXiv:2502.14802v2, 2025.
\newblock \url{https://arxiv.org/abs/2502.14802v2}

\bibitem{arxiv:2410.15944v1}
Ayman Asad Khan, Md Toufique Hasan, Kai Kristian Kemell, Jussi Rasku, Pekka Abrahamsson.
\newblock \emph{Developing Retrieval Augmented Generation (RAG) based LLM Systems from PDFs: An Experience Report}.
\newblock arXiv:2410.15944v1, 2024.
\newblock \url{https://arxiv.org/abs/2410.15944v1}

\end{thebibliography}

\end{document}
